\documentclass[11pt]{article}

% ============================================================
% Packages
% ============================================================
\usepackage[margin=1in]{geometry}
\usepackage{amsmath,amssymb,amsthm,mathtools}
\usepackage{bm}
\usepackage{xcolor}
\usepackage{algorithm}
\usepackage{algpseudocode}
\usepackage[numbers,sort&compress]{natbib}
\usepackage{hyperref}
\usepackage{booktabs}
\usepackage{aliascnt}
\usepackage[nameinlink,capitalize]{cleveref}

\hypersetup{
    colorlinks=true,
    linkcolor=blue!60!black,
    citecolor=blue!60!black,
    urlcolor=blue!60!black,
    pdftitle={Covariance Bootstrap, Quadratic Rescue, and Gaussian-Core Sharp Regions for Gaussian Natural Gradient Variational Inference},
    pdfauthor={}
}

% ============================================================
% Theorem environments
% ============================================================
\theoremstyle{plain}
\newtheorem{theorem}{Theorem}[section]

\newaliascnt{proposition}{theorem}
\newtheorem{proposition}[proposition]{Proposition}
\aliascntresetthe{proposition}

\newaliascnt{lemma}{theorem}
\newtheorem{lemma}[lemma]{Lemma}
\aliascntresetthe{lemma}

\newaliascnt{corollary}{theorem}
\newtheorem{corollary}[corollary]{Corollary}
\aliascntresetthe{corollary}

\theoremstyle{definition}
\newaliascnt{assumption}{theorem}
\newtheorem{assumption}[assumption]{Assumption}
\aliascntresetthe{assumption}

\newaliascnt{definition}{theorem}
\newtheorem{definition}[definition]{Definition}
\aliascntresetthe{definition}

\theoremstyle{remark}
\newaliascnt{remark}{theorem}
\newtheorem{remark}[remark]{Remark}
\aliascntresetthe{remark}

\crefname{theorem}{Theorem}{Theorems}
\Crefname{theorem}{Theorem}{Theorems}
\crefname{proposition}{Proposition}{Propositions}
\Crefname{proposition}{Proposition}{Propositions}
\crefname{lemma}{Lemma}{Lemmas}
\Crefname{lemma}{Lemma}{Lemmas}
\crefname{corollary}{Corollary}{Corollaries}
\Crefname{corollary}{Corollary}{Corollaries}
\crefname{assumption}{Assumption}{Assumptions}
\Crefname{assumption}{Assumption}{Assumptions}
\crefname{definition}{Definition}{Definitions}
\Crefname{definition}{Definition}{Definitions}
\crefname{remark}{Remark}{Remarks}
\Crefname{remark}{Remark}{Remarks}
\crefname{algorithm}{Algorithm}{Algorithms}
\Crefname{algorithm}{Algorithm}{Algorithms}
\crefname{section}{Section}{Sections}
\Crefname{section}{Section}{Sections}
\crefname{appendix}{Appendix}{Appendices}
\Crefname{appendix}{Appendix}{Appendices}

% ============================================================
% Macros
% ============================================================
\newcommand{\R}{\mathbb{R}}
\newcommand{\N}{\mathcal{N}}
\newcommand{\E}{\mathbb{E}}
\newcommand{\Prob}{\mathbb{P}}
\newcommand{\Var}{\operatorname{Var}}
\newcommand{\KL}{\mathrm{KL}}
\newcommand{\Tr}{\operatorname{Tr}}
\newcommand{\SPD}{\mathbb{S}_{++}}
\newcommand{\Sym}{\mathbb{S}}
\newcommand{\Id}{I}
\newcommand{\dd}{\mathrm{d}}
\newcommand{\grad}{\operatorname{grad}}
\newcommand{\op}{\mathrm{op}}
\newcommand{\Frob}{\mathrm{F}}
\newcommand{\post}{\mathrm{post}}
\newcommand{\FR}{\mathrm{FR}}
\newcommand{\BW}{\mathrm{BW}}
\newcommand{\Riem}{\mathrm{R}}
\newcommand{\calE}{\mathcal{E}}
\newcommand{\PG}{\mathcal{P}_{\mathrm G}}
\newcommand{\Aspace}{\mathcal{A}}
\newcommand{\Fil}{\mathcal{F}}
\newcommand{\STL}{\mathrm{STL}}
\newcommand{\norm}[1]{\left\lVert #1\right\rVert}
\newcommand{\inner}[2]{\left\langle #1,#2\right\rangle}
\newcommand{\logp}{\log_+}
\newcommand{\glob}{\mathrm{glob}}
\newcommand{\cov}{\mathrm{cov}}
\newcommand{\loc}{\mathrm{loc}}
\newcommand{\Gradnorm}{\mathsf{G}}
\newcommand{\LH}{L_H}
\newcommand{\barLH}{\overline{L}_H}
\newcommand{\LB}{L_\Lambda}
\newcommand{\Retr}{\operatorname{Retr}}
\newcommand{\Kresc}{\mathsf{K}}
\newcommand{\Fresc}{\mathsf{F}}
\newcommand{\diag}{\operatorname{diag}}
\newcommand{\Ureg}{\mathcal{U}}
\newcommand{\Hreg}{\mathcal{H}}
\newcommand{\FGauss}{F_{\mathrm G}}
\newcommand{\gstar}{\gamma_\star}
\newcommand{\gund}{\underline{\gamma}}
\newcommand{\Lstar}{L_\star}
\newcommand{\sharpthr}{\Delta^{\sharp}}
\newcommand{\Kng}{K_{\mathrm{ng}}}
\newcommand{\Kdisc}{K^{\mathrm{disc}}}
\newcommand{\LHs}{L_{H,\star}}
\newcommand{\PROOFHOLE}{{\color{red}\textsf{[\,proof to be inserted\,]}}}
\newcommand{\TODO}[1]{{\color{red}$\blacktriangleright$~\textsf{TODO:}~\textit{#1}}}
\newcommand{\SKEL}[1]{{\color{blue!55!black}\small$\ast$~\textsf{#1}}}
\newcommand{\alphastar}{\alpha_\star}
\newcommand{\betastar}{\beta_\star}
\newcommand{\kappastar}{\kappa_\star}

% ============================================================
% Title
% ============================================================
\title{Gaussian Variational inference with Fisher--Rao gradient flow}
\author{}
\date{}

\begin{document}
\maketitle

\begin{abstract}
    \TODO \\Abstract
\end{abstract}

\tableofcontents


\section{Introduction}
\label{sec:intro}

\TODO{Opening: target $\pi\propto e^{-V}$ on $\R^{N_\theta}$, the Gaussian family
$\mathcal{Q}=\{\N(m,C)\}$, the KL objective, and Gaussian VI as the intrinsic Fisher--Rao /
affine-invariant natural-gradient flow. Motivate affine invariance as the organizing symmetry and
foreshadow that Bures--Wasserstein enters only as a warm-start device.}

\subsection{Two global obstructions and the global-to-local picture}
\label{sec:intro-picture}
\TODO{(1) Underdispersed covariance $\Rightarrow$ a covariance burn-in phase (residual
$\varepsilon=1-cA$). (2) Intrinsic nonquadratic conditioning obstructs a dimension-free rate even
after whitening. These produce the three-stage decomposition---covariance burn-in, global
localization, local convergence---that organizes the paper; point forward to
\Cref{thm:det-three-stage} (deterministic) and \Cref{cor:stoch-three-stage} (stochastic).}

\subsection{Contributions}
\label{sec:intro-contrib}

\begin{description}
\item[\textnormal{(I) Sharp deterministic global theory.}]
An affine-invariant continuous-time rate in the \emph{intrinsic} condition number $\kappastar$
(\Cref{thm:cont-global-star}), with an order-optimal covariance burn-in and a matching continuous-time
lower bound (\Cref{thm:spiral-sharp,cor:cont-sharp}); and sharp fixed-step discrete complexity for both
the Riemannian retraction and KL/Bregman schemes (\Cref{thm:Riem-global,thm:KL-global}) with a matching
$\Omega(\kappa/\Delta t)$ lower bound (\Cref{thm:disc-global-sharp,cor:disc-global-sharp}). Covariance
burn-in is removed by a Bures--Wasserstein warm start or a single quadratic-rescue Hessian query
(\Cref{thm:disc-W-to-FR,thm:gauss-recovery,thm:disc-rescue-to-FR}), giving the three-stage complexity
(\Cref{thm:det-three-stage}).
\item[\textnormal{(II) Gaussian-core local theory.}]
A decomposition $F=\FGauss+H$ that is exact for the Gaussian part and makes the certified local region
\emph{sharp} for Gaussian targets (\Cref{thm:core-flow,cor:gauss-sharp}), with a discrete local
contraction for both schemes (\Cref{thm:disc-local}).
\item[\textnormal{(III) Stochastic algorithms, end to end.}]
Price/Hessian sticking-the-landing algorithms (\Cref{alg:unified}) with pathwise covariance
confinement; constant-step global convergence in two regimes (\Cref{thm:Riem-selfbound,thm:Riem-Hlip}
and their KL analogues \Cref{thm:KL-selfbound,thm:KL-Hlip}), floor-free decreasing-step convergence
(\Cref{thm:decreasing}), and a stochastic Gaussian-core local contraction with an end-to-end survival
guarantee (\Cref{thm:stoch-local-Riem,thm:stoch-local-KL,cor:stoch-three-stage}).
\end{description}

\subsection{Main results at a glance}
\label{sec:intro-tables}
\TODO{One or two sentences on the reading convention: $\widetilde O/\widetilde\Omega$ hide logarithmic
factors; $\kappastar$ is the affine-invariant intrinsic condition number and $\kappa=\beta/\alpha$ the
extrinsic one, with $1\le\kappastar\le\kappa^2$.}

\begin{table}[ht]
\centering
\renewcommand{\arraystretch}{1.25}
\begin{tabular}{@{}llll@{}}
\toprule
\textbf{Dynamics / scheme} & \textbf{Global complexity} & \textbf{Sharpness} & \textbf{Local behavior}\\
\midrule
Continuous Fisher--Rao flow & $\widetilde O(\kappastar)$ & $\widetilde\Omega(\kappastar)$ (\Cref{thm:spiral-sharp}) & Gaussian-core rate\\
Riemannian retraction (step $\Delta t$) & $\widetilde O(\kappa/\Delta t)$ & $\Omega(\kappa/\Delta t)$ (\Cref{thm:disc-global-sharp}) & geometric\\
KL/Bregman resolvent (step $\Delta t$) & $\widetilde O(\kappa/\Delta t)$ & $\Omega(\kappa/\Delta t)$ (\Cref{thm:disc-global-sharp}) & geometric\\
\bottomrule
\end{tabular}
\caption{Deterministic complexity summary. Global counts are to reach $F-F_\star\le\delta$; local rates
hold on the certified region of \Cref{sec:local}. \TODO{finalize constants / footnote the
input-radius dependence of the lower bounds.}}
\label{tab:det-intro}
\end{table}

\begin{table}[ht]
\centering
\renewcommand{\arraystretch}{1.25}
\begin{tabular}{@{}lllll@{}}
\toprule
\textbf{Regime} & \textbf{Outer iters} & \textbf{Batch $B$} & \textbf{Floor} & \textbf{Reference}\\
\midrule
Self-bounding (SLC only)    & $\widetilde O(\kappa^3)$ & $1$              & $O(\kappa N_\theta)$              & \Cref{thm:Riem-selfbound}\\
Hessian--Lipschitz          & $\widetilde O(\kappa^2)$ & $1$              & $O(\kappa N_\theta^2\barLH^2/B)$  & \Cref{thm:Riem-Hlip}\\
Decreasing step             & $O(1/\delta)$            & $1$              & $0$                               & \Cref{thm:decreasing}\\
Local (non-exit)            & $O((4+\Gamma)^2)$        & $\gtrsim V_1$    & $O(N_\theta^2\LHs^2/B)$           & \Cref{thm:stoch-local-Riem}\\
Gaussian target             & $1$                      & $1$              & $0$                               & \Cref{thm:gauss-recovery,prop:gauss-nofloor}\\
\bottomrule
\end{tabular}
\caption{Stochastic complexity summary for the Price/Hessian STL scheme
($\Gamma=O(1+N_\theta\log\kappa)$, $V_1=24+6N_\theta\LHs^2$); single-sample counts unless a batch is
indicated. \TODO{cross-check floors against \Cref{sec:stoch}.}}
\label{tab:stoch-intro}
\end{table}

\subsection{Related work}
\label{sec:related}
\TODO{Gaussian VI / natural gradient \citep{amari1998natural,opper2009variational}; Bures--Wasserstein
gradient flow and forward--backward \citep{lambert2022variational,diao2023forward}; strongly-log-concave
sampling rates \citep{chen2023sampling}; sticking-the-landing and its variance
\citep{roeder2017sticking,kim2024stl,domke2023provable}. Two brief remarks on KLS-adjacent literature
(not needed under strong log-concavity).}

\subsection{Notation and organization}
\label{sec:org}
\TODO{Notation table; then a one-paragraph roadmap: \Cref{sec:setup} geometry/objective/algorithms;
\Cref{sec:cont-conv} continuous-time theory and its sharpness; \Cref{sec:disc} discretization, global
complexity, and burn-in removal; \Cref{sec:local} the Gaussian-core local theory; \Cref{sec:stoch} the
stochastic theory; \Cref{sec:numerics} experiments; \Cref{sec:discussion} discussion.}


\section{Fisher--Rao geometry, the variational objective, and algorithms}
\label{sec:setup}

We write the energy gap
\[
    \Delta(a):=\calE(a)-\calE(a_\star),
    \qquad
    \Delta_0:=\Delta(a_0),
    \qquad
    \Delta_n:=\Delta(a_n),
\]
and denote by $\Delta t>0$ the Fisher--Rao stepsize (written $h$ in one-step lemmas). Throughout,
$\Delta t$ carries the differential $t$, so there is no clash with the energy gap $\Delta(\cdot)$;
$N_\theta$ is the parameter dimension.

\begin{assumption}[Strong log-concavity and smoothness]
\label{assump:logconcave-smooth}
There exist $0<\alpha\leq\beta<\infty$ with
\[
    \alpha\Id\preceq\nabla^2V(\theta)\preceq\beta\Id,
    \qquad
    \forall\theta\in\R^{N_\theta}.
\]
We write $\kappa:=\beta/\alpha\geq1$.
\end{assumption}

\Cref{assump:logconcave-smooth} is the standing regularity assumption throughout. The local phase
of \Cref{sec:local} additionally uses whitened bounds $\alpha_\star,\beta_\star$ defined there, and
the sharpened stochastic theory of \Cref{sec:stoch-Hlip} additionally uses a Hessian-Lipschitz
assumption (\Cref{assump:global-Hess-Lip}).

For $a=(m,C)\in\Aspace$ we write
\begin{equation}
\label{eq:gAR}
    g(a):=-\E_{\rho_a}[\nabla V(X)],
    \qquad
    A(a):=\E_{\rho_a}[\nabla^2V(X)],
    \qquad
    R(a):=C^{-1}-A(a),
\end{equation}
and call $R(a)$ the \emph{precision residual}. \Cref{assump:logconcave-smooth} gives
$\alpha\Id\preceq A(a)\preceq\beta\Id$. At the minimizer,
\begin{equation}
\label{eq:first-order}
    g(a_\star)=0,
    \qquad
    A(a_\star)=C_\star^{-1},
    \qquad\text{i.e.}\qquad
    R(a_\star)=0 .
\end{equation}
We use the energy splitting $\calE=\calE_1+\calE_2$,
\begin{equation}
\label{eq:splitting}
    \calE_1(a)=\E_{\rho_a}[\log\rho_a]=-\tfrac12\log\det C+\text{const.},
    \qquad
    \calE_2(a)=-\E_{\rho_a}[\log\rho_{\post}]=\E_{\rho_a}[V]+\log Z .
\end{equation}

\subsection{Initial covariance and the band constant}

We assume $\lambda_{0,\min}\Id\preceq C_0\preceq\lambda_{0,\max}\Id$ and set
\begin{equation}
\label{eq:Lambda}
    \lambda_0:=\lambda_{0,\min},
    \qquad
    \Lambda:=\max\Bigl\{\lambda_{0,\max},\frac1\alpha\Bigr\},
    \qquad
    \LB:=\beta\Lambda,
    \qquad
    \logp x:=\max\{\log x,0\}.
\end{equation}
Since $\Lambda\geq1/\alpha$ we always have $\LB\geq\kappa\geq1$.

\subsection{The Gaussian Fisher--Rao metric}
\label{sec:metric}

\begin{definition}[Gaussian Fisher--Rao metric]
\label{def:FR-metric}
For $a=(m,C)\in\Aspace$ and tangent vectors
$(u,X),(v,Y)\in\R^{N_\theta}\times\Sym(N_\theta)$, the intrinsic (information-geometric)
Fisher--Rao metric on $\PG$ is
\begin{equation}
\label{eq:FR-metric}
    g_a\bigl((u,X),(v,Y)\bigr)
    :=
    u^\top C^{-1}v+\frac12\Tr\bigl(C^{-1}XC^{-1}Y\bigr),
\end{equation}
with induced tangent norm
\begin{equation}
\label{eq:tangent-norm}
    \norm{(u,X)}_{a,\FR}^2
    :=
    u^\top C^{-1}u+\frac12\norm{C^{-1/2}XC^{-1/2}}_{\Frob}^2 .
\end{equation}
We abbreviate $\inner{\cdot}{\cdot}_a:=g_a(\cdot,\cdot)$ and $\norm{\cdot}_a:=\norm{\cdot}_{a,\FR}$.
\end{definition}

The factor $\tfrac12$ on the covariance block is the one induced by the Fisher information of the
Gaussian family.

\begin{lemma}[Gaussian covariance derivative; Price's theorem, directional form]
\label{lem:price-derivative}
Let $C\in\SPD(N_\theta)$, $Y\in\Sym(N_\theta)$, and let $V\in C^2$ satisfy
$\norm{\nabla^2V}_\op\leq\beta$. Write $a_\epsilon:=(m,C+\epsilon Y)$. Then
\begin{equation}
\label{eq:price-directional}
    \frac{\dd}{\dd\epsilon}\bigg|_{\epsilon=0}\E_{\rho_{a_\epsilon}}[V]
    =\frac12\Tr\bigl(Y\,A(a)\bigr),
    \qquad
    \frac{\dd}{\dd\epsilon}\bigg|_{\epsilon=0}\calE_1(a_\epsilon)
    =-\frac12\Tr\bigl(YC^{-1}\bigr),
\end{equation}
and consequently
\begin{equation}
\label{eq:E-directional}
    \frac{\dd}{\dd\epsilon}\bigg|_{\epsilon=0}\calE(a_\epsilon)
    =\frac12\Tr\Bigl(Y\bigl(A(a)-C^{-1}\bigr)\Bigr)
    =-\frac12\Tr\bigl(Y R(a)\bigr).
\end{equation}
\end{lemma}

\begin{proof}
\PROOFHOLE
\end{proof}

\begin{lemma}[Fisher--Rao gradient, flow, and gradient norm]
\label{lem:FR-gradient}
Under \Cref{def:FR-metric},
\begin{equation}
\label{eq:FR-grad}
    \grad\calE(a)=\bigl(-Cg(a),\,-CR(a)C\bigr),
\end{equation}
so that the Fisher--Rao (natural) gradient flow $\dot a_t=-\grad\calE(a_t)$ is
\begin{equation}
\label{eq:NG-flow}
    \frac{\dd m_t}{\dd t}=C_tg(a_t),
    \qquad
    \frac{\dd C_t}{\dd t}=C_t-C_tA(a_t)C_t .
\end{equation}
Moreover, defining the Fisher--Rao gradient norm
\begin{equation}
\label{eq:G-norm}
    \Gradnorm(a)^2:=g(a)^\top Cg(a)+\frac12\norm{C^{1/2}R(a)C^{1/2}}_{\Frob}^2 ,
\end{equation}
one has $\norm{\grad\calE(a)}_a^2=\Gradnorm(a)^2$ and, along \eqref{eq:NG-flow},
$\tfrac{\dd}{\dd t}\Delta(a_t)=-\Gradnorm(a_t)^2$.
\end{lemma}

\begin{proof}
\PROOFHOLE
\end{proof}

\subsection{Deterministic schemes}

The deterministic \emph{Riemannian retraction scheme} (covariance-exponential scheme) is
\begin{equation}
\label{eq:Riem-det}
    m_{n+1}=m_n+\Delta t\,C_ng(a_n),
    \qquad
    C_{n+1}=C_n^{1/2}\exp\bigl(\Delta t(\Id-C_n^{1/2}A(a_n)C_n^{1/2})\bigr)C_n^{1/2},
\end{equation}
and the deterministic \emph{KL/Bregman scheme} is
\begin{equation}
\label{eq:KL-det}
    m_{n+1}=m_n+\Delta t\,C_ng(a_n),
    \qquad
    C_{n+1}=(1+\Delta t)\bigl(C_n^{-1}+\Delta t\,A(a_n)\bigr)^{-1},
\end{equation}
equivalently $C_{n+1}^{-1}=\tfrac1{1+\Delta t}\bigl(C_n^{-1}+\Delta t\,A(a_n)\bigr)$.

\begin{remark}[Terminology: retraction, not geodesic]
\label{rem:not-geodesic}
The covariance block $C\mapsto C^{1/2}\exp(\cdot)C^{1/2}$ is the exponential map of the
affine-invariant metric on $\SPD(N_\theta)$, but \eqref{eq:Riem-det} pairs it with a Euclidean
step in the mean, and we do not claim that the resulting map is the exponential map of the full
Gaussian Fisher--Rao metric \eqref{eq:FR-metric}. All statements are proved for the retraction
\eqref{eq:retraction} that the algorithms use; we therefore say \emph{Riemannian retraction
scheme} rather than ``geodesic scheme''.
\end{remark}

\section{Continuous-time convergence and its sharpness}
\label{sec:cont-conv}

\subsection{Two basic inequalities}
\label{sec:basic}

\begin{lemma}[Bures--Wasserstein PL inequality on Gaussians]
\label{lem:W2-PL}
Under \Cref{assump:logconcave-smooth}, for every $a\in\Aspace$,
\[
    \norm{\grad^{W_2}\calE(a)}_{W_2}^2
    =\norm{g(a)}^2+\Tr\bigl(R(a)CR(a)\bigr)
    \geq2\alpha\,\Delta(a),
\]
where $\grad^{W_2}\calE(a)=(-g(a),-R(a))$ is the Gaussian-restricted Bures--Wasserstein gradient.
\end{lemma}

This is the Polyak--{\L}ojasiewicz inequality for the restricted gradient, following from the
$\alpha$-geodesic convexity of $\calE$ along Bures--Wasserstein geodesics between Gaussians; see
\citet{lambert2022variational,diao2023forward}.

\begin{lemma}[Fisher--Rao/Wasserstein comparison]
\label{lem:FR-W2}
Let $a=(m,C)$ with $C\succeq\ell\Id$, $\ell>0$. Then
\begin{equation}
\label{eq:FR-W2}
    \boxed{\;
    \Gradnorm(a)^2
    \;\geq\;
    \frac{\ell}{2}\norm{\grad^{W_2}\calE(a)}_{W_2}^2
    \;\geq\;
    \alpha\ell\,\Delta(a).
    \;}
\end{equation}
Moreover $\Gradnorm(a)^2\leq\norm{C}_\op\norm{\grad^{W_2}\calE(a)}_{W_2}^2$.
\end{lemma}

\begin{proof}
\PROOFHOLE
\end{proof}

\subsection{Continuous-time covariance bootstrap}
\label{sec:cont}

\begin{lemma}[Precision Duhamel formula]
\label{lem:precision-duhamel}
Let $P_t:=C_t^{-1}$. Along \eqref{eq:NG-flow},
$\dot P_t=A(a_t)-P_t$, hence $P_t=e^{-t}P_0+\int_0^te^{-(t-s)}A(a_s)\,\dd s$.
\end{lemma}

\begin{proof}
\PROOFHOLE
\end{proof}

\begin{lemma}[Continuous covariance bootstrap]
\label{lem:cont-cov}
Under \Cref{assump:logconcave-smooth}, along \eqref{eq:NG-flow},
\begin{equation}
\label{eq:ell-t}
    C_t\succeq\ell(t)\Id,
    \qquad
    \ell(t)=\frac{\lambda_0e^t}{1+\beta\lambda_0(e^t-1)},
    \qquad\text{and}\qquad
    C_t\preceq\Lambda\Id .
\end{equation}
In particular $C_t\succeq\tfrac1{2\beta}\Id$ for all $t\geq\logp\frac1{\beta\lambda_0}$.
\end{lemma}

\begin{proof}
\PROOFHOLE
\end{proof}

\begin{theorem}[Continuous-time global convergence; deterministic]
\label{thm:cont-global}
Under \Cref{assump:logconcave-smooth}, the flow \eqref{eq:NG-flow} satisfies
\[
    \boxed{\;
    \Delta(a_t)\leq\Delta_0\bigl(1+\beta\lambda_0(e^t-1)\bigr)^{-1/\kappa}.
    \;}
\]
Consequently, for $\delta\in(0,\Delta_0)$, the hitting time
$T_{\mathrm c}^{\glob}(\delta):=\inf\{t\geq0:\Delta(a_t)\leq\delta\}$ obeys
\[
    T_{\mathrm c}^{\glob}(\delta)
    \leq\log\!\left(1+\frac{(\Delta_0/\delta)^{\kappa}-1}{\beta\lambda_0}\right)
    =O\!\left(\logp\frac1{\beta\lambda_0}+\kappa\logp\frac{\Delta_0}{\delta}\right).
\]
\end{theorem}

\begin{proof}
\PROOFHOLE
\end{proof}

\subsection{Affine-invariant sharpening and continuous-time sharpness}
\label{sec:cont-sharp}

\subsubsection{Optimizer whitening and intrinsic conditioning}

Let $a_\star=(m_\star,C_\star)$ be the unique optimizing Gaussian and set
\begin{equation}
\label{eq:optimizer-whitening}
    \widehat m:=C_\star^{-1/2}(m-m_\star),\qquad
    \widehat C:=C_\star^{-1/2}CC_\star^{-1/2},\qquad
    \widehat V(z):=V(m_\star+C_\star^{1/2}z).
\end{equation}
The optimizer becomes $(\widehat m_\star,\widehat C_\star)=(0,\Id)$. Define the optimal whitened
curvature constants
\begin{equation}
\label{eq:star-curvature}
    \alphastar:=\inf_{z\in\R^{N_\theta}}\lambda_{\min}(\nabla^2\widehat V(z)),\qquad
    \betastar:=\sup_{z\in\R^{N_\theta}}\lambda_{\max}(\nabla^2\widehat V(z)),\qquad
    \kappastar:=\frac{\betastar}{\alphastar},
\end{equation}
and the whitened initial covariance floor
$\lambda_{0,\star}:=\lambda_{\min}(\widehat C_0)=\lambda_{\min}(C_\star^{-1/2}C_0C_\star^{-1/2})$.

\begin{lemma}[Optimizer whitening, flow covariance, and affine invariance]
\label{lem:optimizer-whitening}
Under \Cref{assump:logconcave-smooth}, the transformed parameters solve the same Gaussian Fisher--Rao
system,
\begin{equation}
\label{eq:whitened-flow}
    \dot{\widehat m}_t=\widehat C_t\widehat g(\widehat a_t),\qquad
    \dot{\widehat C}_t=\widehat C_t-\widehat C_t\widehat A(\widehat a_t)\widehat C_t,
\end{equation}
where $\widehat g=-\E[\nabla\widehat V]$ and $\widehat A=\E[\nabla^2\widehat V]$. Moreover
$\E_{Z\sim\N(0,\Id)}[\nabla\widehat V(Z)]=0$ and $\E_{Z\sim\N(0,\Id)}[\nabla^2\widehat V(Z)]=\Id$, so
$\alphastar\leq1\leq\betastar$. The quantities $\alphastar,\betastar,\kappastar$ and the spectrum of
$\widehat C_0$ are invariant under arbitrary invertible affine changes of the original coordinates.
\end{lemma}
\begin{proof}
\PROOFHOLE
\end{proof}

\begin{lemma}[Relation between intrinsic and original condition numbers]
\label{lem:kappa-star-kappa}
Under \Cref{assump:logconcave-smooth}, $\tfrac1\beta\Id\preceq C_\star\preceq\tfrac1\alpha\Id$ and
\begin{equation}
\label{eq:kappa-star-kappa}
    \tfrac1\kappa\Id\preceq\nabla^2\widehat V(z)\preceq\kappa\Id,\qquad 1\leq\kappastar\leq\kappa^2.
\end{equation}
The upper comparison can be very loose: for an anisotropic quadratic target, $\kappastar=1$ even when
$\kappa$ is arbitrarily large.
\end{lemma}
\begin{proof}
\PROOFHOLE
\end{proof}

\begin{lemma}[Whitened covariance floor]
\label{lem:cont-cov-star}
Along \eqref{eq:NG-flow}, in optimizer-whitened coordinates,
$\widehat C_t\succeq\ell_\star(t)\Id$ with
$\ell_\star(t):=\lambda_{0,\star}e^t/\bigl(1+\betastar\lambda_{0,\star}(e^t-1)\bigr)$; in particular
$\widehat C_t\succeq(2\betastar)^{-1}\Id$ for $t\geq\logp\bigl(1/(\betastar\lambda_{0,\star})\bigr)$.
\end{lemma}
\begin{proof}
\PROOFHOLE
\end{proof}

\subsubsection{Affine-invariant global upper bound}

\begin{theorem}[Affine-invariant continuous-time global convergence]
\label{thm:cont-global-star}
Under \Cref{assump:logconcave-smooth}, for every $t\geq0$,
\begin{equation}
\label{eq:cont-star-decay}
    \boxed{\Delta(a_t)\leq\Delta_0\bigl[1+\betastar\lambda_{0,\star}(e^t-1)\bigr]^{-1/\kappastar}.}
\end{equation}
Consequently, for every $\delta\in(0,\Delta_0)$,
\begin{equation}
\label{eq:cont-star-hitting-exact}
    T_{\mathrm c}^{\glob}(\delta):=\inf\{t\geq0:\Delta(a_t)\leq\delta\}
    \leq\log\!\Bigl(1+\tfrac{(\Delta_0/\delta)^{\kappastar}-1}{\betastar\lambda_{0,\star}}\Bigr),
\end{equation}
and in particular
\begin{equation}
\label{eq:cont-star-hitting-simple}
    \boxed{T_{\mathrm c}^{\glob}(\delta)\leq\logp\frac1{\betastar\lambda_{0,\star}}
    +\kappastar\log\frac{\Delta_0}{\delta}.}
\end{equation}
The original-coordinate estimate of \Cref{thm:cont-global} remains valid; one may take the minimum of
the two.
\end{theorem}
\begin{proof}
\PROOFHOLE
\end{proof}

\begin{remark}[What is and is not encoded by $\kappastar$]
The intrinsic condition number removes anisotropy that can be eliminated by an affine
reparametrization. It does not, by itself, encode the initialization size: the factor
$\log(\Delta_0/\delta)$, or an equivalent radius/input-size term, is indispensable, since if arbitrary
initial means are allowed the worst-case time to reach a fixed energy level is infinite even for a
quadratic target.
\end{remark}

\subsubsection{Sharpness of the continuous-time global-localization rate}

\begin{theorem}[Continuous-time logarithmic-spiral lower bound]
\label{thm:spiral-sharp}
For every $K\geq16$ there is a globally $C^3$, even potential $V_K$ on $\R^2$ satisfying
$\alpha\Id\preceq\nabla^2V_K\preceq\beta\Id$ with $\kappa=2K+1$ (the logarithmic-spiral construction,
deferred to \Cref{app:spiral}), together with an admissible anisotropic initialization satisfying the
general spectral condition $(1/\beta)\Id\preceq C_0\preceq(1/\alpha)\Id$, such that for any local level
$\delta_{\mathrm{loc}}<\tfrac\alpha2 e^{-0.08}\norm{m_0}^2$ the exact Gaussian Fisher--Rao flow obeys
\begin{equation}
\label{eq:spiral-lower-time}
    \boxed{T_{\mathrm{loc}}:=\inf\{t\geq0:\Delta(a_t)\leq\delta_{\mathrm{loc}}\}
    \geq\frac K{100}\geq\frac{\kappastar(V_K)}{103}.}
\end{equation}
\end{theorem}
\begin{proof}
\PROOFHOLE
\end{proof}

\begin{corollary}[Sharp polynomial dependence in dimension at least two]
\label{cor:cont-sharp}
For $N_\theta\geq2$, the worst-case polynomial dependence of the continuous-time global-localization
phase on the affine-invariant condition number is
\begin{equation}
\label{eq:cont-sharp-rate}
    \boxed{T_{\mathrm{loc}}=\widetilde\Theta(\kappastar),}
\end{equation}
under the general spectral initialization class, where the tilde suppresses logarithmic dependence on
initialization size and localization accuracy. For the spiral family, $\kappastar=\Theta(\kappa)$ and
$T_{\mathrm{loc}}\geq(\kappa-1)/200$.
\end{corollary}
\begin{proof}
\PROOFHOLE
\end{proof}

\begin{remark}[Scope of the sharpness theorem]
The matching construction uses an admissible anisotropic initial covariance. It does not establish an
$\Omega(\kappastar)$ lower bound under the stricter initialization $C_0=\beta^{-1}\Id$. It also allows
the initial radius to grow with $K$ (and, in a Hessian--Lipschitz variant, with $1/L_H$); the logarithm
of that radius is part of the input size charged by the upper theorem. These restrictions should be
retained whenever the sharpness result is quoted.
\end{remark}

\section{Discretization: global complexity, sharpness, and burn-in removal}
\label{sec:disc}

\subsection{The Riemannian one-step descent lemma}
\label{sec:descent}

Throughout this section $a=(m,C)$ is fixed and
\begin{equation}
\label{eq:S-def}
    S:=\Id-C^{1/2}A(a)C^{1/2}=C^{1/2}R(a)C^{1/2},
\end{equation}
so that, by \eqref{eq:G-norm}, $\Gradnorm(a)^2=g(a)^\top Cg(a)+\tfrac12\norm{S}_{\Frob}^2$. We write
the retraction
\begin{equation}
\label{eq:retraction}
    \Retr_a(u,U):=\Bigl(m+u,\;C^{1/2}\exp\bigl(C^{-1/2}UC^{-1/2}\bigr)C^{1/2}\Bigr),
\end{equation}
which reproduces \eqref{eq:Riem-det} with stepsize $h$ on taking
$(u,U)=-h\,\grad\calE(a)$, that is $u=hCg$ and $U=hCRC$,
since $C^{-1/2}(hCRC)C^{-1/2}=hS$.

\subsubsection{Two scalar inequalities}

\begin{lemma}[Scalar inequality, deterministic step]
\label{lem:scalar-det}
Let $K\geq1$ and $0<h\leq1/K$. Then for every $s\in[1-K,1]$,
\begin{equation}
\label{eq:scalar-det}
    \bigl(e^{hs/2}-1\bigr)(1-s)+\frac K2\bigl(e^{hs/2}-1\bigr)^2-\frac h2s
    \;\leq\;
    -\frac h2\Bigl(1-\frac{Kh}2\Bigr)s^2 .
\end{equation}
\end{lemma}

\begin{proof}
\PROOFHOLE
\end{proof}

\begin{lemma}[Scalar inequality, arbitrary retraction direction]
\label{lem:scalar-stoch}
For every $w\in[-1,1]$,
\begin{equation}
\label{eq:scalar-stoch}
    \Bigl(e^{w/2}-1-\frac w2\Bigr)+\frac12\bigl(e^{w/2}-1\bigr)^2\;\leq\;\frac{w^2}2 .
\end{equation}
\end{lemma}

\begin{proof}
\PROOFHOLE
\end{proof}

\subsubsection{Deterministic one-step descent}

\begin{lemma}[Deterministic Riemannian one-step descent]
\label{lem:det-descent}
Assume $0\preceq\nabla^2V\preceq\beta\Id$ and $C\preceq\Lambda\Id$; set $K:=\beta\Lambda$ and
assume $K\geq1$. Let $0<h\leq1/K$ and
$m^+:=m+hCg(a)$, $C^+:=C^{1/2}\exp(hS)C^{1/2}$, $a^+:=(m^+,C^+)$. Then
\begin{equation}
\label{eq:det-descent}
    \boxed{\;
    \calE(a^+)\;\leq\;\calE(a)-h\Bigl(1-\frac{Kh}2\Bigr)\Gradnorm(a)^2 .
    \;}
\end{equation}
\end{lemma}

\begin{proof}
\PROOFHOLE
\end{proof}

\subsubsection{Retraction smoothness for arbitrary directions}

\begin{lemma}[Retraction smoothness, constant $2\beta\Lambda$]
\label{lem:retr-smooth}
Assume $0\preceq\nabla^2V\preceq\beta\Id$ and $C\preceq\Lambda\Id$; set $K:=\beta\Lambda$ and
$L:=2K$. Let $\zeta=(u,U)$ be any tangent vector at $a$ with
\begin{equation}
\label{eq:scope}
    \norm{W}_\op\leq1,\qquad W:=C^{-1/2}UC^{-1/2}.
\end{equation}
Then
\begin{equation}
\label{eq:retr-smooth}
    \boxed{\;
    \calE\bigl(\Retr_a(\zeta)\bigr)
    \;\leq\;
    \calE(a)+\inner{\grad\calE(a)}{\zeta}_a+\frac L2\norm{\zeta}_a^2 .
    \;}
\end{equation}
\end{lemma}

\begin{proof}
\PROOFHOLE
\end{proof}

\subsection{Riemannian retraction scheme: deterministic global theory}
\label{sec:Riem-global}

\begin{lemma}[Riemannian covariance bootstrap]
\label{lem:Riem-cov}
Assume \Cref{assump:logconcave-smooth}, $\lambda_{0,\min}\Id\preceq C_0\preceq\lambda_{0,\max}\Id$
and $0<\Delta t\leq1/\LB$. Then $C_n\preceq\Lambda\Id$ for all $n\geq0$. Define
\begin{equation}
\label{eq:Riem-envelope}
    L_0:=\min\Bigl\{\lambda_0,\frac1{2\beta}\Bigr\},
    \qquad
    L_{n+1}:=\frac{e^{\Delta t}L_n}{1+(e-1)\Delta t\,\beta L_n}.
\end{equation}
Then $C_n\succeq L_n\Id$ for all $n\geq0$, and with $y_n:=(\beta L_n)^{-1}$,
\begin{equation}
\label{eq:Riem-y}
    y_n=y_\infty+(y_0-y_\infty)e^{-n\Delta t},
    \qquad
    y_\infty:=\frac{(e-1)\Delta t}{e^{\Delta t}-1}\in[1,e-1)\subset[1,2).
\end{equation}
Consequently $C_n\succeq\tfrac1{2\beta}\Id$ for all
\begin{equation}
\label{eq:Ncov-minus}
    n\geq N_{\cov,-}^{\Riem}
    :=\left\lceil\frac1{\Delta t}\logp\frac{y_0-y_\infty}{2-y_\infty}\right\rceil
    =O\!\left(\frac1{\Delta t}\logp\frac1{\beta\lambda_0}\right),
\end{equation}
and, with $z_0:=\tfrac1{\alpha\lambda_{0,\max}}$ and $z_\infty:=\tfrac{\Delta t}{e^{\Delta t}-1}\in[\tfrac1{e-1},1)$,
$C_n\preceq\tfrac2\alpha\Id$ for all $n\geq N_{\cov,+}^{\Riem}=O(1/\Delta t)$, where
$N_{\cov,+}^{\Riem}=0$ if $z_0\geq\tfrac12$ and
$N_{\cov,+}^{\Riem}=\lceil\Delta t^{-1}\log\tfrac{z_\infty-z_0}{z_\infty-1/2}\rceil$ otherwise.
\end{lemma}

\begin{proof}
\PROOFHOLE
\end{proof}

\begin{theorem}[Deterministic Riemannian global contraction]
\label{thm:Riem-global}
Assume \Cref{assump:logconcave-smooth} and $0<\Delta t\leq1/\LB$. Along \eqref{eq:Riem-det}, with
$L_n$ from \eqref{eq:Riem-envelope}:
\begin{enumerate}
\item For every $n\geq0$,
$\Delta_{n+1}\leq\bigl(1-\alpha L_n\Delta t(1-\tfrac{\LB\Delta t}2)\bigr)\Delta_n$.
\item For $n\geq N_{\cov,-}^{\Riem}$ one has $L_n\geq\tfrac1{2\beta}$ and
$\Delta_{n+1}\leq\bigl(1-\tfrac{\Delta t}{4\kappa}\bigr)\Delta_n$.
\item For $\delta\in(0,\Delta_0)$,
$N_{\Riem}^{\glob}(\delta)\leq N_{\cov,-}^{\Riem}+\lceil\tfrac{4\kappa}{\Delta t}\log\tfrac{\Delta_0}\delta\rceil
=O\bigl(\tfrac1{\Delta t}\logp\tfrac1{\beta\lambda_0}+\tfrac\kappa{\Delta t}\logp\tfrac{\Delta_0}\delta\bigr)$.
\end{enumerate}
\end{theorem}

\begin{proof}
\PROOFHOLE
\end{proof}

\subsection{KL/Bregman scheme: deterministic global theory}
\label{sec:KL-global}

\begin{lemma}[KL covariance bootstrap]
\label{lem:KL-cov}
Assume \Cref{assump:logconcave-smooth} and $\lambda_{0,\min}\Id\preceq C_0\preceq\lambda_{0,\max}\Id$.
Along \eqref{eq:KL-det}, $C_n\preceq\Lambda\Id$ for all $n\geq0$. With
$L_0:=\min\{\lambda_0,\tfrac1{2\beta}\}$ and $L_{n+1}=\tfrac{(1+\Delta t)L_n}{1+\Delta t\,\beta L_n}$,
one has $C_n\succeq L_n\Id$ with
$L_n=\bigl(\beta[1+(\tfrac1{\beta L_0}-1)(1+\Delta t)^{-n}]\bigr)^{-1}$. Hence $C_n\succeq\tfrac1{2\beta}\Id$
for all
$n\geq N_{\cov,-}^{\KL}:=\lceil\tfrac{\logp(\tfrac1{\beta L_0}-1)}{\log(1+\Delta t)}\rceil
=O(\tfrac1{\Delta t}\logp\tfrac1{\beta\lambda_0})$,
and $C_n\preceq\tfrac2\alpha\Id$ for all $n\geq N_{\cov,+}^{\KL}:=\lceil\tfrac{\log2}{\log(1+\Delta t)}\rceil=O(1/\Delta t)$.
\end{lemma}

\begin{proof}
\PROOFHOLE
\end{proof}

\begin{proposition}[KL/Bregman one-step estimate]
\label{prop:KL-onestep}
Let $0<\Delta t\leq1/\LB$ and $L_n\Id\preceq C_n\preceq\Lambda\Id$. Then, in the forward divergence,
\[
    \calE(a_{n+1})\leq\calE(a_n)-\frac1{\Delta t}\KL(\rho_{a_{n+1}}\|\rho_{a_n}),
    \qquad
    \KL(\rho_{a_{n+1}}\|\rho_{a_n})
    \geq\frac{L_n\Delta t^2}{4(1+\Delta t\,\LB)^2}\norm{\grad^{W_2}\calE(a_n)}_{W_2}^2 .
\]
\end{proposition}

The KL map \eqref{eq:KL-det} is a \emph{forward--backward} (proximal-gradient) step in the Gaussian
KL/Bregman geometry, not a full-gradient mirror step: it linearizes the potential energy $\calE_2$
explicitly and treats the entropy $\calE_1$ implicitly, which is what produces the implicit resolvent
covariance update $C_{n+1}=(1+\Delta t)(C_n^{-1}+\Delta t\,A(a_n))^{-1}$. \Cref{app:KL-onestep} proves
the proposition in full and self-contained form: the variational characterization and explicit maps
(\Cref{lem:KL-argmin}), a scalar inequality for the covariance factor (\Cref{lem:scalar-KL}), and the
one-step estimate itself (\Cref{lem:KL-onestep-app})---both the descent margin, proved directly by a
Gaussian coupling of $\rho_{a_n}$ and $\rho_{a_{n+1}}$ together with \Cref{lem:scalar-KL}, and the
lower-bound factor $\tfrac{L_n\Delta t^2}{4(1+\Delta t\LB)^2}$. Unlike the $W_2^2$-proximal
Bures--Wasserstein forward--backward scheme of \citet{diao2023forward,lambert2022variational} (whose
transport warm-start is used separately in \Cref{sec:transport-bootstrap}), the descent here is in the
\emph{forward} KL divergence $\KL(\rho_{a_{n+1}}\|\rho_{a_n})$ and the constants are specific to the
$\tfrac12$-Fisher--Rao normalization.

\begin{theorem}[Deterministic KL/Bregman global contraction]
\label{thm:KL-global}
Assume \Cref{assump:logconcave-smooth} and $0<\Delta t\leq1/\LB$. Along \eqref{eq:KL-det}:
\begin{enumerate}
\item For every $n\geq0$,
$\Delta_{n+1}\leq\bigl(1-\tfrac{\alpha L_n\Delta t}{2(1+\Delta t\,\LB)^2}\bigr)\Delta_n$.
\item For $n\geq N_{\cov,-}^{\KL}$, $\Delta_{n+1}\leq\bigl(1-\tfrac{\Delta t}{16\kappa}\bigr)\Delta_n$.
\item $N_{\KL}^{\glob}(\delta)\leq N_{\cov,-}^{\KL}+\lceil\tfrac{16\kappa}{\Delta t}\log\tfrac{\Delta_0}\delta\rceil
=O\bigl(\tfrac1{\Delta t}\logp\tfrac1{\beta\lambda_0}+\tfrac\kappa{\Delta t}\logp\tfrac{\Delta_0}\delta\bigr)$.
\end{enumerate}
\end{theorem}

\begin{proof}
\PROOFHOLE
\end{proof}

\subsection{Sharpness of fixed-step discrete global localization}
\label{sec:disc-sharp}

\begin{theorem}[Sharp fixed-step global localization for both discretizations]
\label{thm:disc-global-sharp}
Fix $\gamma\in(0,1/2]$ and $L_0>0$. There exist constants $T,q>0$ and $\kappa_0\geq1$, depending only on
$\gamma$, $L_0$, and a fixed flat-top curvature bump $\varphi$, such that for every $\kappa\geq\kappa_0$
there is an even potential $V_\kappa\in C^\infty(\R)$ satisfying
\begin{equation}
\label{eq:disc-sharp-final-assumptions}
    1\leq V_\kappa''\leq\kappa,\qquad \|V_\kappa'''\|_\infty\leq L_0,
\end{equation}
and an initialization $m_0=Y\kappa^4/\gamma$, $c_0=1/\kappa$ (the bump-train construction, deferred to
\Cref{app:bumptrain}), such that, when either fixed-step scheme is run with exact expectations and
stepsize $h=\gamma/\kappa$,
\begin{equation}
\label{eq:disc-sharp-final-gap}
    \Delta_n\geq q\,\Delta_0\qquad\text{for every }0\leq n\leq\bigl\lfloor T\kappa^2/\gamma\bigr\rfloor.
\end{equation}
Consequently, for constant-factor objective reduction,
\begin{equation}
\label{eq:disc-sharp-final-complexity}
    N_{\Riem}=\Omega\!\Bigl(\tfrac{\kappa^2}{\gamma}\Bigr)=\Omega\!\Bigl(\tfrac{\kappa}{h}\Bigr),\qquad
    N_{\KL}=\Omega\!\Bigl(\tfrac{\kappa^2}{\gamma}\Bigr)=\Omega\!\Bigl(\tfrac{\kappa}{h}\Bigr).
\end{equation}
The lower bound holds with $\kappastar=\kappa$, with the covariance initially in the curvature-scale
band, and with a dimension-free Hessian--Lipschitz constant.
\end{theorem}
\begin{proof}
\PROOFHOLE
\end{proof}

\begin{corollary}[Sharpness of the global-localization prefactor]
\label{cor:disc-global-sharp}
For the globally fixed-step Riemannian and KL/Bregman schemes, the factor $\kappa/\Delta t$ in
\Cref{thm:Riem-global,thm:KL-global} is sharp, up to constants independent of $\kappa$ for each fixed
$\gamma=\kappa\Delta t$, for a constant-factor reduction of the energy. In particular, with
$\Delta t=\Theta(\kappa^{-1})$, both methods have worst-case global-localization complexity
$\Theta(\kappa^2)$ in the condition-number exponent. This remains true under the additional assumption
$\|V'''\|_\infty=O(1)$.
\end{corollary}
\begin{proof}
\PROOFHOLE
\end{proof}

\begin{remark}[Limitations of the lower bound]
The theorem proves sharpness of the polynomial condition-number dependence for explicit methods
committed to one globally fixed step. It does not rule out an $O(\kappa)$ global-localization result for
an adaptive relative-curvature step, an Armijo or backtracking line search, or an implicit or
Rosenbrock-stabilized mean update: the construction has $c_nA(m_n,c_n)\simeq1$, so a method whose
stability test is based on the current relative curvature could accept an order-one step and avoid this
particular slowdown.
\end{remark}

\subsection{Removing the covariance burn-in: transport and quadratic rescue}
\label{sec:transport-bootstrap}

The Fisher--Rao schemes converge globally from any positive-definite initialization, incurring only
the logarithmic covariance burn-in $\logp\tfrac1{\beta\lambda_0}$ of
\Cref{thm:cont-global,thm:Riem-global,thm:KL-global}. When removing that stage is desirable, two
optional \emph{deterministic} initializers are available, and we develop both here. The first is a
short Bures--Wasserstein transport warm-start (\Cref{subsec:transport-warmstart}): it is
expectation-level, uses only $g(a)$ and $A(a)$, and requires no pointwise Hessian. The second, under
a pointwise Hessian oracle, is the \emph{quadratic rescue} (\Cref{subsec:rescue-det}); this is
precisely the single-query initializer that later drives the stochastic pipeline of
\Cref{sec:stoch-schemes}, and we give it here in its deterministic role. Neither initializer is part
of the mandatory algorithm and neither affects the local discretization factor; each merely deletes
the $\logp\tfrac1{\beta\lambda_0}$ stage. They differ in oracle model and in geometry, as recorded in
\Cref{rem:two-initializers}.

\subsubsection{Bures--Wasserstein transport warm-start}
\label{subsec:transport-warmstart}

For $\dot C=C-CAC$, an eigenvalue $\mu_i$ obeys $\dot\mu_i=\mu_i(1-\mu_ih_i)$, $h_i\in[\alpha,\beta]$;
in the underdispersed regime $\mu_i\ll1/\beta$ this gives $\dot\mu_i\approx\mu_i$ (multiplicative),
so reaching $1/\beta$ from $\lambda_0$ costs time $\asymp\log(1/(\beta\lambda_0))$. The Gaussian
Bures--Wasserstein flow
\begin{equation}
\label{eq:W-flow}
    \dot m_t=g(a_t),\qquad \dot C_t=2\Id-C_tA(a_t)-A(a_t)C_t
\end{equation}
gives $\dot\mu_i=2(1-\mu_ih_i)\approx2$ (additive, scale-independent). Transport is thus efficient
for the covariance warm-up $0\to O(1/\beta)$, Fisher--Rao for shape calibration $1/\beta\to1/\alpha$.
The flow \eqref{eq:W-flow} is not affine invariant, so the scale $1/\beta$ is stated in the given
coordinates.

\begin{lemma}[Wasserstein covariance bootstrap]
\label{lem:W-cov-bootstrap}
Under \Cref{assump:logconcave-smooth}, along \eqref{eq:W-flow},
$C_t\succeq\ell_W(t)\Id$ with $\ell_W(t):=\tfrac1\beta+(\lambda_0-\tfrac1\beta)e^{-2\beta t}$, and
$C_t\preceq\Lambda\Id$. Consequently, for $c\in(0,1)$,
$T_W^{\mathrm c}(c):=\tfrac1{2\beta}\log\tfrac{1-\beta\lambda_0}{1-c}$ if $\lambda_0<c/\beta$ (and $0$
otherwise) satisfies $C_{T_W^{\mathrm c}(c)}\succeq\tfrac c\beta\Id$; the warm-up time is $O(1/\beta)$,
with no logarithmic dependence on $\lambda_0$.
\end{lemma}

\begin{proof}
\PROOFHOLE
\end{proof}

\begin{lemma}[One-step Wasserstein covariance floor]
\label{lem:one-step-W}
Given $a=(m,C)$, $\eta>0$, put $M:=\Id-\eta A(a)$, $m^+:=m+\eta g(a)$, $\widetilde C:=MCM$, and
\begin{equation}
\label{eq:W-FB-cov}
    C^+:=\tfrac12\bigl(\widetilde C+2\eta\Id+[\widetilde C(\widetilde C+4\eta\Id)]^{1/2}\bigr).
\end{equation}
If $C\preceq\Lambda\Id$ and $0<\eta\leq1/\beta$, then $C^+\succeq\eta\Id$, $C^+\preceq\Lambda\Id$, and
the Bures--Wasserstein forward--backward descent inequality gives $\calE(m^+,C^+)\leq\calE(m,C)$.
\end{lemma}

\begin{proof}
\PROOFHOLE
\end{proof}

\begin{theorem}[Discrete transport-warm-started natural gradient; deterministic]
\label{thm:disc-W-to-FR}
Assume \Cref{assump:logconcave-smooth}. Fix $c\in(0,1]$, $\eta:=c/\beta$. If $\lambda_0<c/\beta$
perform one step \eqref{eq:W-FB-cov}; otherwise skip it. Then run \eqref{eq:Riem-det} with
$0<\Delta t\leq\min\{1/\LB,\,2/(4+\Gamma),\,1/(\beta_\star\bar\Lambda_{\delta_{\Riem}})\}$. For small
$\epsilon>0$,
\[
    \boxed{
    N_{W\to\Riem}(\epsilon;c)
    =\mathbf 1_{\{\lambda_0<c/\beta\}}
    +O\!\left(\frac\kappa{c\,\Delta t}\logp\frac{\Delta_0}{\delta_{\Riem}}+\frac{4+\Gamma}{\Delta t}\log\frac1\epsilon\right),
    }
\]
with $\delta_{\Riem}$ a level admissible for \Cref{thm:disc-local}. The same holds for
\eqref{eq:KL-det} with $\delta_{\Riem}\to\delta_{\KL}$ and $4+\Gamma\to4+2\Delta t+\Gamma$.
\end{theorem}

\begin{proof}
\PROOFHOLE
\end{proof}

\begin{remark}
\label{rem:transport-caveat}
The transport warm-start removes the small-initial-covariance burn-in. It does not improve the
local discretization factor $(4+\Gamma)/\Delta t$; the local stiffness is separate from the global
covariance warm-up.
\end{remark}

\subsubsection{Deterministic quadratic rescue}
\label{subsec:rescue-det}

The second burn-in remover is available whenever a pointwise Hessian can be queried at a single
point. It is exactly the initializer used by the stochastic pipeline of \Cref{sec:stoch-schemes};
we state and analyse it here in its deterministic role.

\begin{definition}[Quadratic rescue]
\label{def:rescue}
Given $m_{\mathrm{in}}\in\R^{N_\theta}$, query the oracle at $m_{\mathrm{in}}$:
$G_{\mathrm{in}}:=\nabla V(m_{\mathrm{in}})$, $H_{\mathrm{in}}:=\nabla^2V(m_{\mathrm{in}})$, and set
\begin{equation}
\label{eq:rescue}
    \boxed{\;
    m_0:=m_{\mathrm{in}}-H_{\mathrm{in}}^{-1}G_{\mathrm{in}},
    \qquad
    C_0:=H_{\mathrm{in}}^{-1}.
    \;}
\end{equation}
\end{definition}

\begin{lemma}[Curvature-scale covariance rescue]
\label{lem:rescue-band}
Under \Cref{assump:logconcave-smooth}, \eqref{eq:rescue} gives $\tfrac1\beta\Id\preceq C_0\preceq\tfrac1\alpha\Id$.
\end{lemma}

\begin{proof}
\PROOFHOLE
\end{proof}

\begin{theorem}[Exact one-query recovery of Gaussian targets]
\label{thm:gauss-recovery}
Suppose $V(x)=\tfrac12(x-\mu)^\top H(x-\mu)+c$ with $H\succ0$. Then \eqref{eq:rescue} gives $m_0=\mu$,
$C_0=H^{-1}$, so $a_0=a_\star$ and $\Delta(a_0)=0$, for both the Riemannian and KL schemes. Over the
class of unknown Gaussian targets, the oracle complexity is exactly one.
\end{theorem}

\begin{proof}
\PROOFHOLE
\end{proof}

\begin{theorem}[Quadratic-rescued deterministic natural gradient; no covariance burn-in]
\label{thm:disc-rescue-to-FR}
Assume \Cref{assump:logconcave-smooth}, and initialize by the quadratic rescue of
\Cref{def:rescue}. By \Cref{lem:rescue-band} then $\tfrac1\beta\Id\preceq C_0\preceq\tfrac1\alpha\Id$,
so $\lambda_0:=\lambda_{0,\min}\geq\tfrac1\beta$. Consequently, in the covariance envelopes of
\Cref{lem:Riem-cov,lem:KL-cov}, $L_0:=\min\{\lambda_0,\tfrac1{2\beta}\}=\tfrac1{2\beta}$ and
$y_0:=(\beta L_0)^{-1}=2$, so the covariance floor $\tfrac1{2\beta}\Id$ holds from $n=0$ and the
burn-in count is $N_{\mathrm{cov},-}^{\Riem}=N_{\mathrm{cov},-}^{\KL}=0$. Running \eqref{eq:Riem-det}
(resp.\ \eqref{eq:KL-det}) with
$0<\Delta t\leq\min\{1/\LB,\,2/(4+\Gamma),\,1/(\beta_\star\bar\Lambda_{\delta_\bullet})\}$, for small
$\epsilon>0$,
\[
    \boxed{\;
    N_{\mathrm{rescue}\to\bullet}(\epsilon)
    =1
    +O\!\left(\frac\kappa{\Delta t}\logp\frac{\Delta_0}{\delta_\bullet}
    +\frac{c_\bullet}{\Delta t}\log\frac1\epsilon\right),
    \qquad
    c_{\Riem}=4+\Gamma,\quad c_{\KL}=4+2\Delta t+\Gamma,
    \;}
\]
where the additive $1$ counts the single deterministic Hessian query, $\delta_\bullet$ is a level
admissible for \Cref{thm:disc-local}, and $\Delta_0:=\Delta(a_0)$ is finite (bounded in
\Cref{lem:rescue-energy} under \Cref{assump:global-Hess-Lip}, and $\Delta_0=0$ for Gaussian targets
by \Cref{thm:gauss-recovery}). In particular the $\logp\tfrac1{\beta\lambda_0}$ burn-in stage of
\Cref{thm:det-three-stage} is eliminated.
\end{theorem}

\begin{proof}
\PROOFHOLE
\end{proof}

\begin{remark}[Two initializers, one oracle distinction]
\label{rem:two-initializers}
Both warm-starts delete the $\logp\tfrac1{\beta\lambda_0}$ stage but sit in different oracle models
and geometries. The transport warm-start (\Cref{thm:disc-W-to-FR}) is \emph{expectation-level}: it
uses $g(a)$ and $A(a)=\E_{\rho_a}[\nabla^2V]$ and inflates the covariance additively along the
Bures--Wasserstein forward--backward step, so it never queries a pointwise Hessian and only guarantees
a covariance \emph{floor} $C_0\succeq\tfrac c\beta\Id$ (the lower band edge; $c=1$ gives $\tfrac1\beta\Id$),
without matching the pointwise curvature. The quadratic rescue (\Cref{thm:disc-rescue-to-FR}) is a single
\emph{pointwise-Hessian} (local-curvature/Newton) query and lands directly on the full curvature
band $[\tfrac1\beta,\tfrac1\alpha]$, exact for Gaussian targets. Thus transport trades a weaker
oracle for a coarser landing scale, while the rescue trades one Hessian evaluation for exact
Gaussian recovery; the ensuing Fisher--Rao iterations are identical.
\end{remark}

\subsection{Deterministic three-stage complexity}
\label{sec:three-stage}

\begin{theorem}[Deterministic three-stage complexity]
\label{thm:det-three-stage}
Assume \Cref{assump:logconcave-smooth} and $\lambda_{0,\min}\Id\preceq C_0\preceq\lambda_{0,\max}\Id$.
Let $\delta_{\mathrm c}>0$ be admissible for \Cref{thm:core-flow} with
$\rho_{\mathrm{core}}(\delta_{\mathrm c})\geq\gund$, and $\delta_{\Riem},\delta_{\KL}>0$ admissible for
\Cref{thm:disc-local}; such levels exist by \Cref{cor:core-universal} and \Cref{thm:disc-local}.
For sufficiently small $\epsilon>0$:
\begin{enumerate}
\item \textup{(Continuous flow.)}
$T_{\mathrm c}(\epsilon)=O\bigl(\logp\tfrac1{\beta\lambda_0}+\kappa\logp\tfrac{\Delta_0}{\delta_{\mathrm c}}+(4+\Gamma)\log\tfrac1\epsilon\bigr)$
gives $\norm{a_T-a_\star}_\star^2\leq\epsilon$.
\item \textup{(Riemannian, $0<\Delta t\leq\min\{1/\LB,2/(4+\Gamma),1/(\beta_\star\bar\Lambda_{\delta_{\Riem}})\}$.)}
$N_{\Riem}(\epsilon)=O\bigl(\tfrac1{\Delta t}\logp\tfrac1{\beta\lambda_0}+\tfrac\kappa{\Delta t}\logp\tfrac{\Delta_0}{\delta_{\Riem}}+\tfrac{4+\Gamma}{\Delta t}\log\tfrac1\epsilon\bigr)$.
\item \textup{(KL, same window.)}
$N_{\KL}(\epsilon)=O\bigl(\tfrac1{\Delta t}\logp\tfrac1{\beta\lambda_0}+\tfrac\kappa{\Delta t}\logp\tfrac{\Delta_0}{\delta_{\KL}}+\tfrac{4+2\Delta t+\Gamma}{\Delta t}\log\tfrac1\epsilon\bigr)$.
\end{enumerate}
\end{theorem}

\begin{proof}
\PROOFHOLE
\end{proof}

\begin{corollary}[Burn-in-free three-stage complexity]
\label{cor:burnin-free}
Under \Cref{assump:logconcave-smooth}, initialize by either the quadratic rescue
(\Cref{thm:disc-rescue-to-FR}) or the transport warm-start (\Cref{thm:disc-W-to-FR}) \emph{with target
$c=1$} (i.e.\ $\eta=1/\beta$, at most one forward--backward step). Both land $C_0\succeq\tfrac1\beta\Id$,
so $\lambda_0\geq\tfrac1\beta$, and the first stage of \Cref{thm:det-three-stage} is removed: for both
discretizations
\[
    N_\bullet(\epsilon)=n_{\mathrm{init}}
    +O\!\left(\frac\kappa{\Delta t}\logp\frac{\Delta_0}{\delta_\bullet}
    +\frac{c_\bullet}{\Delta t}\log\frac1\epsilon\right),
    \qquad c_{\Riem}=4+\Gamma,\quad c_{\KL}=4+2\Delta t+\Gamma,
\]
where $n_{\mathrm{init}}\in\{0,1\}$ is the warm-start ($1$ Hessian query for the rescue, or
$\mathbf 1_{\{\lambda_0<1/\beta\}}\leq1$ forward--backward steps for transport) and $\Delta_0=\Delta(a_0)$
is the energy of the warm-started initializer; for Gaussian targets the rescue gives $\Delta_0=0$ and
the localization stage vanishes as well.
\end{corollary}

\begin{proof}
\PROOFHOLE
\end{proof}

\section{Local convergence via the Gaussian core}
\label{sec:local}

The organizing principle of the local theory is a \emph{Gaussian-core decomposition}. The natural
gradient field of the standard Gaussian target is nonlinear in the covariance but can be dissipated
exactly; treating it as a Taylor remainder discards a large, genuinely dissipative contribution. We
write
\[
    F=\FGauss+H,
    \qquad
    \FGauss(m,C):=(-Cm,\;C-C^2),
\]
retain $\FGauss$ exactly, and Taylor-bound only the non-Gaussian perturbation $H$. Because
$H\equiv0$ whenever the target is Gaussian, the certified energy-sublevel region is sharp for
Gaussian targets.

\subsection{Whitened coordinates and the linearized gap}
\label{subsec:whitened}

Define $\widehat m:=C_\star^{-1/2}(m-m_\star)$, $\widehat C:=C_\star^{-1/2}CC_\star^{-1/2}$, and
$\widehat V(z):=V(m_\star+C_\star^{1/2}z)$. Let $0<\alpha_\star\leq\beta_\star<\infty$ satisfy
$\alpha_\star\Id\preceq\nabla^2\widehat V(z)\preceq\beta_\star\Id$, with
$\kappa_\star:=\beta_\star/\alpha_\star$. By \eqref{eq:first-order},
$\E_{\N(0,\Id)}[\nabla\widehat V]=0$ and $\E_{\N(0,\Id)}[\nabla^2\widehat V]=\Id$, so
$\alpha_\star\leq1\leq\beta_\star$, and one may take $\alpha_\star\geq1/\kappa$,
$\beta_\star\leq\kappa$, $\kappa_\star\leq\kappa^2$. Throughout this section we work in whitened
coordinates: $a=(m,C)$ denotes the whitened parameter, $a_\star=(0,\Id)$. Statements transfer back
through the fixed equilibrium norms
\begin{equation}
\label{eq:local-norms}
    \norm{a-a_\star}_\star^2:=\norm{\widehat m}^2+\tfrac12\norm{\widehat C-\Id}_{\Frob}^2,
    \qquad
    \norm{a-a_\star}_{\star,\Delta t}^2:=\norm{\widehat m}^2+\tfrac{1+\Delta t}2\norm{\widehat C-\Id}_{\Frob}^2 .
\end{equation}
We write $\xi:=a-a_\star$, $D(a):=\norm{\xi}_\star^2$, and $F(a):=(Cg(a),\,C-CA(a)C)$, so
\eqref{eq:NG-flow} reads $\dot a=F(a)$, $F(a_\star)=0$. Set
\begin{equation}
\label{eq:Gamma}
    \Gamma:=\min\Bigl\{\beta_\star-1,\,N_\theta\bigl(3+\tfrac4{\sqrt\pi}(1+\log\kappa_\star)\bigr)\Bigr\},
    \qquad
    \gund:=\frac1{4+\Gamma} .
\end{equation}

\begin{definition}[Linearized score operators]
\label{def:linear-score-ops}
Let $Z\sim\N(0,\Id)$ and $B(Z):=\nabla^2\widehat V(Z)$, so $\E B(Z)=\Id$. For $u\in\R^{N_\theta}$,
$X\in\Sym(N_\theta)$ define
$\mathcal T[X]_k:=\E[\Tr(XB(Z))Z_k]$, $\mathcal T^*[u]_{ij}:=\E[(u^\top Z)B_{ij}(Z)]$, and
$\mathcal H[X]_{ij}:=\E[B_{ij}(Z)(Z^\top XZ-\Tr X)]$. Gaussian score identities give
$u^\top\mathcal T[X]=\Tr(\mathcal T^*[u]X)$ and $\Tr(\mathcal H[X]Y)=\Tr(X\mathcal H[Y])$.
\end{definition}

\begin{lemma}[Diagonal-mode bounds]
\label{lem:linear-diagonal-bounds}
Fix an orthogonal basis in which $X=\diag(\lambda)$, and set $\widetilde T_{ki}:=\E[B_{ii}(Z)Z_k]$,
$G_{ij}:=\E[B_{ii}(Z)Z_j^2]$, $D:=G-\mathbf1\mathbf1^\top$. Then $D$ is symmetric,
\begin{equation}
\label{eq:T-Schur}
    \widetilde T^\top\widetilde T\preceq\Id+D,
    \qquad\text{and}\qquad
    \lambda_{\max}(D)\leq\Gamma .
\end{equation}
\end{lemma}

\begin{proof}
\PROOFHOLE
\end{proof}

\begin{proposition}[Linearized gap]
\label{prop:linearized}
Let $J_\star:=DF(a_\star)$ and $\Lstar:=-J_\star$. Then
\[
    \Lstar(u,X)=\Bigl(u+\tfrac12\mathcal T[X],\;X+\mathcal T^*[u]+\tfrac12\mathcal H[X]\Bigr),
\]
$\Lstar$ is self-adjoint in $\inner{\cdot}{\cdot}_\star$, and with
$\gstar:=\lambda_{\min}(\Lstar)$, $\Lambda_\star:=\lambda_{\max}(\Lstar)$,
\begin{equation}
\label{eq:linear-gap-bounds}
    \boxed{\;\frac1{4+\Gamma}\leq\gstar\leq\Lambda_\star\leq\frac{4+\Gamma}2 .\;}
\end{equation}
For a Gaussian target, $\mathcal T=0$, $\mathcal H=0$, and $\gstar=\Lambda_\star=1$.
\end{proposition}

\begin{proof}
\PROOFHOLE
\end{proof}

\begin{proposition}[Linearized contractions of the deterministic maps]
\label{prop:linearized-maps}
Let $T_{\Riem,h}$, $T_{\KL,h}$ be the maps \eqref{eq:Riem-det}, \eqref{eq:KL-det} in whitened
coordinates. If $0<h\leq2/(4+\Gamma)$, then
\[
    \|DT_{\Riem,h}(a_\star)\xi\|_\star\leq\Bigl(1-\tfrac h{4+\Gamma}\Bigr)\|\xi\|_\star,
    \qquad
    \|DT_{\KL,h}(a_\star)\xi\|_{\star,h}\leq\Bigl(1-\tfrac h{4+2h+\Gamma}\Bigr)\|\xi\|_{\star,h}.
\]
\end{proposition}

\begin{proof}
\PROOFHOLE
\end{proof}

\subsection{Energy geometry: entropy-enhanced coercivity and bands}

\begin{lemma}[Entropy-enhanced coercivity]
\label{lem:entropy-coerc}
Write $B:=C^{1/2}$ with eigenvalues $s_1,\dots,s_{N_\theta}$ and
$\phi_{\alpha_\star}(s):=s-1-\log s+\tfrac{\alpha_\star}2(s-1)^2$. Then
\begin{equation}
\label{eq:entropy-coerc}
    \Delta(m,C)\geq\frac{\alpha_\star}2\norm{m}^2+\sum_{i=1}^{N_\theta}\phi_{\alpha_\star}(s_i).
\end{equation}
For a Gaussian target ($\alpha_\star=1$) this is an equality and
$\Delta_{\mathrm G}(m,C)=\tfrac12\norm{m}^2+\sum_i\varphi(c_i)$ with $\varphi(c):=\tfrac12(c-1-\log c)$
and $c_i$ the eigenvalues of $C$.
\end{lemma}

\begin{proof}
\PROOFHOLE
\end{proof}

\begin{corollary}[Covariance band and norm conversion]
\label{cor:band}
$\phi_{\alpha_\star}$ is strictly decreasing on $(0,1)$, increasing on $(1,\infty)$, with minimum
$0$ at $1$. For $\delta>0$ let $0<s_-(\delta)<1<s_+(\delta)$ solve $\phi_{\alpha_\star}(s)=\delta$,
and set $\ell_\delta:=s_-(\delta)^2$, $U_\delta:=s_+(\delta)^2$. Then $\Delta(a)\leq\delta$ implies
$\ell_\delta\Id\preceq C\preceq U_\delta\Id$. With
\[
    \vartheta(\delta):=\sup_{s\in[s_-,s_+],\,s\neq1}\frac{(s^2-1)^2}{2\phi_{\alpha_\star}(s)}
    \leq\frac{(1+s_+)^2}{\alpha_\star+s_+^{-2}},
    \qquad
    \mathfrak b(\delta)^2:=\max\Bigl\{\tfrac2{\alpha_\star},\vartheta(\delta)\Bigr\},
\]
one has $\Delta(a)\leq\delta\Rightarrow\norm{a-a_\star}_\star\leq\mathfrak b(\delta)\sqrt\delta$, and
$\norm{a-a_\star}_{\star,\Delta t}\leq\mathfrak b_{\Delta t}(\delta)\sqrt\delta$ with
$\mathfrak b_{\Delta t}(\delta)^2:=\max\{2/\alpha_\star,(1+\Delta t)\vartheta(\delta)\}$. For a
Gaussian target, $\ell_\delta,U_\delta$ are the roots of $\varphi(c)=\delta$.
\end{corollary}

\begin{proof}
\PROOFHOLE
\end{proof}

\begin{definition}[Regions, hull, moduli]
\label{def:region}
For $\delta>0$, $\Ureg_\delta:=\{a:\Delta(a)\leq\delta\}$ and
$\Hreg_\delta:=\{a_\star+s(a-a_\star):a\in\Ureg_\delta,s\in[0,1]\}$. Set
$\bar\Lambda_\delta:=\max\{U_\delta,1/\alpha_\star\}$ and the non-Gaussian moduli
\[
    \mathcal B_\delta:=\sup_{a\in\Hreg_\delta}\norm{A(a)-\Id}_{\op},
    \quad
    \Sigma_\delta:=\sup_{a\in\Hreg_\delta}\bigl(\E_{\rho_a}\norm{\nabla^2\widehat V(X)-A(a)}_{\Frob}^2\bigr)^{1/2},
    \quad
    \bar\Sigma_\delta:=\min\{\Sigma_\delta,\beta_\star-\alpha_\star\}.
\]
All three vanish for Gaussian targets and are nondecreasing in $\delta$.
\end{definition}

\begin{lemma}[Geometry and invariance]
\label{lem:region-geom}
Let $\delta>0$, $a=(m,C)\in\Ureg_\delta$. Then (i) $\ell_\delta\Id\preceq C\preceq U_\delta\Id$ and
$\norm{a-a_\star}_\star\leq\mathfrak b(\delta)\sqrt\delta$; (ii) every $a'\in\Hreg_\delta$ satisfies
the same bounds; (iii) $\Ureg_\delta$ is compact; (iv) $\Ureg_\delta$ is forward invariant for the
flow, and for any scheme step with $\calE(a_{n+1})\leq\calE(a_n)$.
\end{lemma}

\begin{proof}
\PROOFHOLE
\end{proof}

\subsection{The exact Gaussian core}

\begin{lemma}[Gaussian-core dissipation]
\label{lem:gauss-core}
Let $\FGauss(m,C):=(-Cm,\,C-C^2)$, the natural gradient field of the standard Gaussian target
$\widehat V_{\mathrm G}(z)=\tfrac12\norm z^2$. Then:
\begin{enumerate}
\item $D\FGauss(a_\star)=-\Id$ on $\R^{N_\theta}\times\Sym(N_\theta)$.
\item For every $a=(m,C)$, with $\xi=a-a_\star$,
$\inner{\xi}{\FGauss(a)}_\star=-m^\top Cm-\tfrac12\Tr\bigl(C(C-\Id)^2\bigr)=-\tfrac12q_{\mathrm G}(a)D(a)$,
where $q_{\mathrm G}(a):=\dfrac{2m^\top Cm+\Tr(C(C-\Id)^2)}{\norm m^2+\tfrac12\norm{C-\Id}_{\Frob}^2}$.
\item $q_{\mathrm G}(a)\geq2\lambda_{\min}(C)$, with equality when $m=0$ and $C-\Id$ is supported on
a minimal eigendirection of $C$; hence on $\Ureg_\delta$,
$\inner{\xi}{\FGauss(a)}_\star\leq-\ell_\delta D(a)$.
\end{enumerate}
\end{lemma}

\begin{proof}
\PROOFHOLE
\end{proof}

\subsection{The non-Gaussian perturbation}

Write $H:=F-\FGauss$. Since $g_{\mathrm G}(a)=-m$ and $A_{\mathrm G}(a)=\Id$,
\begin{equation}
\label{eq:H-def}
    H(a)=\bigl(C(g(a)+m),\;-C(A(a)-\Id)C\bigr),\qquad H(a_\star)=0 .
\end{equation}
Both components vanish identically when $\widehat V=\widehat V_{\mathrm G}$.

\begin{lemma}[Score identities and the sharp second-score constant]
\label{lem:second-score}
Fix $a=(m,C)$, $\zeta=(u,U)$, $v:=C^{-1/2}u$, $S:=C^{-1/2}UC^{-1/2}$,
$t:=\norm\zeta_a=(\norm v^2+\tfrac12\norm S_{\Frob}^2)^{1/2}$, $X=m+C^{1/2}Z$. With
$q_\zeta:=v^\top Z+\tfrac12(Z^\top SZ-\Tr S)$ and
$\chi_\zeta:=q_\zeta^2+(\tfrac12\Tr S^2-\norm v^2-2v^\top SZ-Z^\top S^2Z)$, for $\phi$ of polynomial
growth and $\Phi(a):=\E_{\rho_a}[\phi(X)]$,
\[
    D\Phi(a)[\zeta]=\E[\phi\,q_\zeta],
    \qquad
    D^2\Phi(a)[\zeta,\zeta]=\E[\phi\,\chi_\zeta].
\]
Moreover $q_\zeta$ has Hermite chaoses $1,2$ and $\chi_\zeta$ has chaoses $2,3,4$; in particular
$\E[q_\zeta]=\E[\chi_\zeta]=0$, $\E[Z\chi_\zeta]=0$, $\E[q_\zeta^2]=t^2$, and
$\norm{\chi_\zeta}_{L^2}\leq\sqrt6\,t^2$ (sharp, attained by rank-one covariance directions).
\end{lemma}

\begin{proof}
\PROOFHOLE
\end{proof}

\begin{lemma}[Non-Gaussian second-differential modulus]
\label{lem:Kng}
Define
\[
    \Kng(\delta):=\sup_{a\in\Hreg_\delta}\;\sup_{\zeta\neq0}\;\frac{\norm{D^2H(a)[\zeta,\zeta]}_\star}{\norm\zeta_\star^2}.
\]
Then
\[
    \Kng(\delta)\leq\bigl(K_{\mathrm{ng},m}(\delta)^2+\tfrac12K_{\mathrm{ng},C}(\delta)^2\bigr)^{1/2},
\]
with
\[
    K_{\mathrm{ng},m}(\delta):=\frac{U_\delta^{3/2}}{\ell_\delta^2}\bigl(\sqrt3\bar\Sigma_\delta+\Sigma_\delta+\sqrt{\Sigma_\delta^2+2\mathcal B_\delta^2}\bigr),
    \quad
    K_{\mathrm{ng},C}(\delta):=\frac{U_\delta^2}{\ell_\delta^2}\bigl((4\sqrt2+\sqrt6)\Sigma_\delta+4\mathcal B_\delta\bigr).
\]
In particular $\Kng(\delta)=0$ for Gaussian targets, and, with the spectral-band factors
$U_\delta,\ell_\delta$ of \Cref{cor:band} held fixed,
$\Kng(\delta)=O\!\bigl(\tfrac{U_\delta^2}{\ell_\delta^2}(\mathcal B_\delta+\Sigma_\delta)\bigr)$; the
band factors must be retained in the $O(\cdot)$ since they degrade as $\delta$ grows.
\end{lemma}

\begin{proof}
\PROOFHOLE
\end{proof}

\begin{corollary}[Non-Gaussian Taylor remainder]
\label{cor:taylor-ng}
For $a=a_\star+\xi\in\Ureg_\delta$, write $F(a)=J_\star\xi+N_{\mathrm G}(\xi)+R_{\mathrm{ng}}(\xi)$
with $N_{\mathrm G}(\xi):=\FGauss(a)-J_{\mathrm G}\xi$ ($J_{\mathrm G}=-\Id$) and
$R_{\mathrm{ng}}(\xi):=H(a)-DH(a_\star)\xi$. Then
$\norm{R_{\mathrm{ng}}(\xi)}_\star\leq\tfrac12\Kng(\delta)\norm\xi_\star^2$.
\end{corollary}

\begin{proof}
\PROOFHOLE
\end{proof}

\subsection{Main continuous-time local theorem and sharpness}

\begin{theorem}[Gaussian-core local convergence of the flow]
\label{thm:core-flow}
Let $0<\gamma_0\leq\gstar$ and $\delta>0$, and set
\begin{equation}
\label{eq:rho-core}
    \rho_{\mathrm{core}}(\delta):=2\gamma_0-2+2\ell_\delta-\Kng(\delta)\,\mathfrak b(\delta)\sqrt\delta .
\end{equation}
If $\rho_{\mathrm{core}}(\delta)>0$ and $\Delta(a_0)\leq\delta$, then the flow \eqref{eq:NG-flow}
remains in $\Ureg_\delta$ and $D(a_t)\leq e^{-\rho_{\mathrm{core}}(\delta)t}D(a_0)$ for all $t\geq0$.
\end{theorem}

\begin{proof}
\PROOFHOLE
\end{proof}

\begin{corollary}[Universal certified level]
\label{cor:core-universal}
With $\gamma_0=\gund$, $\rho_{\mathrm{core}}(\delta)\to2\gund$ as $\delta\downarrow0$, so for every
$\rho<2\gund$ there is $\delta_{\mathrm c}(\rho)>0$ with $\rho_{\mathrm{core}}(\delta)\geq\rho$ on
$(0,\delta_{\mathrm c}(\rho)]$.
\end{corollary}

\begin{corollary}[Exact sharpness for Gaussian targets]
\label{cor:gauss-sharp}
Let the target be Gaussian. Then $\gstar=1$, $\Kng\equiv0$, and \Cref{thm:core-flow} holds with the
exact rate $\rho_{\mathrm{core}}(\delta)=2\ell_\delta$. Consequently, for every $\rho\in(0,2)$, the
largest energy sublevel on which rate $\rho$ is certified is $\{\Delta\leq\sharpthr_{\mathrm G}(\rho)\}$
with
\begin{equation}
\label{eq:gauss-sharp-thresh}
    \sharpthr_{\mathrm G}(\rho)=\varphi\Bigl(\tfrac\rho2\Bigr)=\tfrac12\Bigl(\tfrac\rho2-1-\log\tfrac\rho2\Bigr),
\end{equation}
independent of dimension, and this threshold is attained (\Cref{prop:sharp}).
\end{corollary}

\begin{proof}
\PROOFHOLE
\end{proof}

\begin{definition}[Instantaneous rate and maximal threshold]
\label{def:sharp}
For $a\neq a_\star$, set
\[
    q_{\mathrm c}(a):=\frac{-2\inner{a-a_\star}{F(a)}_\star}{\norm{a-a_\star}_\star^2},
    \qquad
    \sharpthr(\rho):=\inf\{\Delta(a):a\neq a_\star,\;q_{\mathrm c}(a)<\rho\}.
\]
\end{definition}

\begin{proposition}[The benchmark is exact and dominates the certified level]
\label{prop:sharp}
Let $\rho>0$. (1) If $\Delta(a_0)<\sharpthr(\rho)$ then $D(a_t)\leq e^{-\rho t}D(a_0)$ for all
$t\geq0$. (2) For every $\delta>\sharpthr(\rho)$ there is $a_0$ with $\Delta(a_0)\leq\delta$ and
instantaneous rate $<\rho$; so $\sharpthr(\rho)$ is the largest energy threshold guaranteeing rate
$\rho$. (3) For a Gaussian target, $\sharpthr(\rho)=\sharpthr_{\mathrm G}(\rho)=\varphi(\rho/2)$ for
$\rho\in(0,2)$, in every dimension, and the certified level of \Cref{cor:gauss-sharp} attains it.
\end{proposition}

\begin{proof}
\PROOFHOLE
\end{proof}

\begin{remark}[Rate versus region]
\label{rem:rate-region}
For Gaussian targets both region and rate are exact. The universal bound $\gund=1/(4+\Gamma)$ is
used only for general targets as a lower bound on $\gstar$; when the exact gap is available it
should be used in \eqref{eq:rho-core}. As $\delta\downarrow0$, $\rho_{\mathrm{core}}(\delta)\to2\gstar$,
the exact asymptotic squared-norm rate.
\end{remark}

\subsection{Deterministic discrete schemes: exact Gaussian intervals}
\label{subsec:disc-local}

\begin{proposition}[Exact Gaussian KL map and sharp interval]
\label{prop:kl-gauss}
Let the target be Gaussian, so \eqref{eq:KL-det} becomes
$m_{n+1}=(\Id-\Delta t\,C_n)m_n$, $C_{n+1}=(1+\Delta t)C_n(\Id+\Delta t\,C_n)^{-1}$. In the
eigenbasis of $C_n$ the scalar factors are
$r_m(c)=(1-\Delta t\,c)^2$, $r_C(c)=(1+\Delta t\,c)^{-2}$, $c^+=\tfrac{(1+\Delta t)c}{1+\Delta t\,c}$.
Fix $0<\Delta t<2$ and $q\in(0,1)$ with $q\geq q_0(\Delta t):=\max\{(1+\Delta t)^{-2},(1-\Delta t)^2\}$,
and set $I_{\Delta t}^{\mathrm{KL}}(q):=[\tfrac{q^{-1/2}-1}{\Delta t},\tfrac{1+\sqrt q}{\Delta t}]\ni1$.
Then $r_m(c),r_C(c)\leq q$ on $I_{\Delta t}^{\mathrm{KL}}(q)$; $c\mapsto c^+$ maps
$I_{\Delta t}^{\mathrm{KL}}(q)$ into itself; and if $\operatorname{spec}(C_n)\subseteq I_{\Delta t}^{\mathrm{KL}}(q)$
then $D_{\Delta t}(a_{n+k})\leq q^kD_{\Delta t}(a_n)$ for all $k\geq0$. For $\Delta t\leq\sqrt2$,
$q_0(\Delta t)=(1+\Delta t)^{-2}$.
\end{proposition}

\begin{proof}
\PROOFHOLE
\end{proof}

\begin{proposition}[General-target discrete local convergence]
\label{thm:disc-local}
Fix a scheme $\bullet\in\{\Riem,\KL\}$ with certified linearized gap $\gamma_\bullet$
($\gamma_{\Riem}=\gund$, $\gamma_{\KL}=\tfrac1{4+2\Delta t+\Gamma}$) and norm $\norm\cdot_{\star,\bullet}$
($\norm\cdot_\star$ resp.\ $\norm\cdot_{\star,\Delta t}$), with associated norm-conversion constant
$\mathfrak b_{\star,\bullet}(\delta)$ defined by $\Delta(a)\leq\delta\Rightarrow\norm{a-a_\star}_{\star,\bullet}\leq\mathfrak b_{\star,\bullet}(\delta)\sqrt\delta$:
explicitly $\mathfrak b_{\star,\Riem}(\delta):=\mathfrak b(\delta)$ and
$\mathfrak b_{\star,\KL}(\delta):=\mathfrak b_{\Delta t}(\delta)$, both of \Cref{cor:band}, with the clean
comparison $\mathfrak b_{\Delta t}(\delta)\leq\sqrt{1+\Delta t}\,\mathfrak b(\delta)$ (immediate from
$\mathfrak b_{\Delta t}(\delta)^2=\max\{2/\alpha_\star,(1+\Delta t)\vartheta(\delta)\}\leq(1+\Delta t)\mathfrak b(\delta)^2$).
For the scheme map $T$ put
$\Kdisc(\delta):=\Delta t^{-1}\sup_{a\in\Hreg_\delta}\norm{D^2T(a)}_{\star,\bullet}$, finite by
\Cref{rem:Kdisc-finite}, where $\norm{D^2T(a)}_{\star,\bullet}:=\sup_{\norm\xi_{\star,\bullet}=1}\norm{D^2T(a)[\xi,\xi]}_{\star,\bullet}$
is the operator norm of the bilinear map $D^2T(a)$ in the scheme norm. Suppose
\[
    0<\Delta t\leq\min\Bigl\{\tfrac2{4+\Gamma},\tfrac1{\beta_\star\bar\Lambda_\delta}\Bigr\},
    \qquad
    \Kdisc(\delta)\,\mathfrak b_{\star,\bullet}(\delta)\sqrt\delta\leq\gamma_\bullet .
\]
If $\Delta(a_0)\leq\delta$, then the iterates remain in $\Ureg_\delta$ and
$\norm{a_n-a_\star}_{\star,\bullet}\leq(1-\tfrac12\Delta t\,\gamma_\bullet)^n\norm{a_0-a_\star}_{\star,\bullet}$.
Since $\Kdisc(\delta)\to\Kdisc(0)<\infty$ and $\mathfrak b_{\star,\bullet}(\delta)\sqrt\delta\to0$, the
radius condition holds for all small $\delta$.
\end{proposition}

\begin{proof}
\PROOFHOLE
\end{proof}

\begin{lemma}[One-step local map contraction before exit]
\label{lem:local-map-contraction}
Fix $\bullet\in\{\Riem,\KL\}$ and a level $\delta>0$ satisfying the radius condition
$\Kdisc(\delta)\,\mathfrak b_{\star,\bullet}(\delta)\sqrt\delta\leq\gamma_\bullet$ of \Cref{thm:disc-local}.
If $0<\Delta t\leq\tfrac2{4+\Gamma}$, then for \emph{every} $a\in\Ureg_\delta$ the deterministic scheme
map $T_\bullet$ satisfies
\[
    \norm{T_\bullet(a)-a_\star}_{\star,\bullet}\leq\Bigl(1-\tfrac12\Delta t\,\gamma_\bullet\Bigr)\norm{a-a_\star}_{\star,\bullet}.
\]
No energy-descent hypothesis is used; in particular the bound requires \emph{not}
$\Delta t\leq1/(\beta_\star\bar\Lambda_\delta)$, and no forward-invariance is asserted---this is purely
the one-step contraction of $T_\bullet$ at the single point $a$.
\end{lemma}

\begin{proof}
\PROOFHOLE
\end{proof}

\begin{remark}[Finiteness and continuity of $\Kdisc$]
\label{rem:Kdisc-finite}
The bound $\Kdisc(\delta)<\infty$ requires more than compactness of $\Hreg_\delta$: it needs
$D^2T$ \emph{continuous} on a neighborhood of the hull. Under \Cref{assump:logconcave-smooth} the
Gaussian-expectation maps $g(a)=-\E_{\rho_a}[\nabla V]$ and $A(a)=\E_{\rho_a}[\nabla^2V]$ are obtained
by differentiating $\E_{\rho_a}[\,\cdot\,]$ under the integral, which is licensed for observables of at
most polynomial growth (here $\nabla V$ grows linearly); each parameter derivative brings down a
Gaussian score polynomial, so $g$ and $A$ are $C^\infty$ in $(m,C)$ on every spectrally bounded compact
subset of $\Aspace$. Both scheme maps $T$ (the Riemannian matrix-exponential retraction and the KL
resolvent $C\mapsto(1+\Delta t)(C^{-1}+\Delta t A)^{-1}$) are smooth functions of $g,A,C$ with $C$
bounded away from $0$ and $\infty$ on $\Hreg_\delta$, hence $C^2$ there with $D^2T$ continuous. The
supremum of a continuous function over the compact hull is finite, giving $\Kdisc(\delta)<\infty$;
continuity in $\delta$ and shrinking of the hull as $\delta\downarrow0$ give
$\Kdisc(\delta)\to\Kdisc(0)<\infty$.
\end{remark}

\subsection{Online entry criteria}
\label{subsec:entry}

The energy test $\Delta(a_n)\leq\delta$ is not directly computable, since $\calE(a_\star)$ is
unknown. The following sufficient tests use only the current iterate and known target quantities,
in original coordinates. They are \emph{exact-expectation} tests: they presuppose the population
quantities $g(a)$ and $A(a)$; a Monte Carlo implementation would replace them by estimators and
require a concentration-adjusted (confidence) version of the gate. Moreover the threshold
$\delta_{\mathrm{loc}}$ they compare against is only \emph{implicitly} available in general (it must
first be certified admissible for \Cref{thm:disc-local}, whose constant $\Kdisc(\delta)$ is bounded
but not given in closed form), so the test is executable only after such a $\delta_{\mathrm{loc}}$
has been fixed. The conservative bounds of \Cref{rem:threshold-computable} give a precomputable
\emph{continuous-time} threshold; for the discrete schemes the admissible level remains implicit
until an explicit upper bound on $\Kdisc(\delta)$ is supplied.

\begin{proposition}[Exact-expectation residual gate, original coordinates]
\label{prop:entry-residual}
At $a=(m,C)$ set $g(a)=-\E_{\rho_a}[\nabla V(X)]$, $A(a)=\E_{\rho_a}[\nabla^2V(X)]$, and
$\mathcal R_{\mathrm{BW}}(a)^2:=\norm{g(a)}^2+\Tr[(C^{-1}-A(a))C(C^{-1}-A(a))]$. Then
$\mathcal R_{\mathrm{BW}}(a)^2\geq2\alpha\Delta(a)$, so for any a priori certified local level
$\delta_{\mathrm{loc}}$,
\[
    \mathcal R_{\mathrm{BW}}(a_n)^2\leq2\alpha\,\delta_{\mathrm{loc}}
    \;\Longrightarrow\;
    \Delta(a_n)\leq\delta_{\mathrm{loc}} .
\]
The left-hand side requires neither $C_\star$, $m_\star$, nor $\calE(a_\star)$. When a covariance
floor $C_n\succeq\ell_n\Id$ is available, the Fisher--Rao gradient norm gives an alternative
intrinsic test, $\Gradnorm(a_n)^2\leq\alpha\ell_n\delta_{\mathrm{loc}}\Rightarrow\Delta(a_n)\leq\delta_{\mathrm{loc}}$,
which by \Cref{lem:FR-W2} is more restrictive than the Bures--Wasserstein gate. For the continuous
flow and either deterministic discretization, once either test triggers the corresponding
energy-descent property keeps $\Delta(a_k)\leq\delta_{\mathrm{loc}}$ at all subsequent times or
iterations. For the stochastic schemes, neither pathwise nor expected containment follows from this
gate alone: the one-step recursion $\E[\Delta_{n+1}\mid\Fil_n]\leq q\Delta_n+\eta$ with $\eta>0$ gives
only $\E[\Delta_{n+1}\mid\Fil_n]\leq q\delta_{\mathrm{loc}}+\eta$, which is $\leq\delta_{\mathrm{loc}}$
only if $\delta_{\mathrm{loc}}\geq\eta/(1-q)$, i.e.\ if the threshold dominates the corresponding
stochastic floor. An expected-containment statement therefore requires either that condition or a
separate stochastic local-stability argument.
\end{proposition}

\begin{proof}
\PROOFHOLE
\end{proof}

\begin{remark}[Precomputable conservative threshold]
\label{rem:threshold-computable}
The level $\delta_{\mathrm{loc}}$ is target-adaptive through
$\alpha_\star,\beta_\star,\Sigma_\delta,\mathcal B_\delta,\Kng$. A precomputable conservative
choice follows from the a priori bounds $\alpha_\star\geq\kappa^{-1}$, $\beta_\star\leq\kappa$,
$\Sigma_\delta\leq\sqrt{N_\theta}(\beta_\star-\alpha_\star)$,
$\mathcal B_\delta\leq\max\{\beta_\star-1,1-\alpha_\star\}$, which reduce the \emph{continuous-time}
admissibility condition $\rho_{\mathrm{core}}(\delta)>0$ of \Cref{thm:core-flow} to a scalar
inequality in $\alpha,\beta,N_\theta$. This does \emph{not} make the \emph{discrete} admissibility
condition $\Kdisc(\delta)\,\mathfrak b_{\star,\bullet}(\delta)\sqrt\delta\leq\gamma_\bullet$ explicit,
because no closed-form upper bound on $\Kdisc(\delta)=\Delta t^{-1}\sup_{a\in\Hreg_\delta}\norm{D^2T(a)}_{\star,\bullet}$
is supplied here; an explicit bound on $\Kdisc(\delta)$ in terms of the target parameters is still
required for a fully precomputable discrete switching threshold.
\end{remark}

\section{Stochastic algorithms}
\label{sec:stoch}

\subsection{The Price/Hessian oracle and the stochastic schemes}
\label{sec:stoch-schemes}

\begin{definition}[Price/Hessian oracle]
\label{def:oracle}
Given $X\in\R^{N_\theta}$, the oracle returns the pair $\bigl(\nabla V(X),\,\nabla^2V(X)\bigr)$.
\end{definition}

Fix a batch size $B\geq1$. Let $\Fil_n:=\sigma(a_0,\ldots,a_n)$, $\xi_{n,s}\overset{\text{iid}}{\sim}\N(0,\Id)$,
$X_{n,s}:=m_n+C_n^{1/2}\xi_{n,s}$, $H_{n,s}:=\nabla^2V(X_{n,s})$.

\subparagraph{Mean/score estimator: STL.}
\begin{equation}
\label{eq:STL-score}
    \overline b_n:=\frac1B\sum_{s=1}^B\bigl(-\nabla V(X_{n,s})+C_n^{-1}(X_{n,s}-m_n)\bigr)
    =\frac1B\sum_{s=1}^B\bigl(-\nabla V(X_{n,s})+C_n^{-1/2}\xi_{n,s}\bigr).
\end{equation}

\subparagraph{Covariance estimator: Price/Hessian.}
$\overline H_n:=\tfrac1B\sum_{s=1}^BH_{n,s}$, $\overline K_n:=C_n^{-1}-\overline H_n$. Both are
conditionally unbiased,
\begin{equation}
\label{eq:STL-unbiased}
    \E[\overline b_n\mid\Fil_n]=g(a_n),
    \qquad
    \E[\overline K_n\mid\Fil_n]=R(a_n).
\end{equation}

\subparagraph{The global update.}
The \emph{minibatch Price/Hessian Riemannian scheme} is
\begin{equation}
\label{eq:Riem-STL}
    m_{n+1}=m_n+\Delta t\,C_n\overline b_n,
    \qquad
    C_{n+1}=C_n^{1/2}\exp\bigl(\Delta t(\Id-C_n^{1/2}\overline H_nC_n^{1/2})\bigr)C_n^{1/2},
\end{equation}
and the \emph{minibatch Price/Hessian KL scheme} is
\begin{equation}
\label{eq:KL-STL}
    m_{n+1}=m_n+\Delta t\,C_n\overline b_n,
    \qquad
    C_{n+1}=(1+\Delta t)\bigl(C_n^{-1}+\Delta t\,\overline H_n\bigr)^{-1}.
\end{equation}
STL enters through the score $\overline b_n$; the covariance part is Price/Hessian. Since each
$H_{n,s}$ has spectrum in $[\alpha,\beta]$ a.s., so does the average:
\begin{equation}
\label{eq:pathwise-Hess}
    \alpha\Id\preceq\overline H_n\preceq\beta\Id\qquad\text{almost surely.}
\end{equation}
This almost-sure spectral bound is the structural reason the Price/Hessian oracle admits an
unconditional global theory.

\begin{remark}[Price's identity and the choice of covariance oracle]
\label{rem:price}
By Gaussian integration by parts, for $X=m+C^{1/2}\xi$, $\E[\xi\,\nabla V(X)^\top]=C^{1/2}A(a)$, so
$A(a)$ admits both a sampled-Hessian and a gradient-only representation. We use the sampled Hessian:
it is the representation with the almost-sure spectral bound \eqref{eq:pathwise-Hess}, which is
exactly what the pathwise confinement arguments require. The gradient-only representation has
unbounded samplewise spectrum; none of the pathwise statements below apply to it.
\end{remark}

\subsubsection{The quadratic rescue}
\label{subsec:rescue}

The stochastic pipeline is initialized by the single deterministic Price/Hessian query of the
\emph{quadratic rescue} (\Cref{def:rescue}), developed for the deterministic theory in
\Cref{subsec:rescue-det}: it sets $m_0=m_{\mathrm{in}}-H_{\mathrm{in}}^{-1}G_{\mathrm{in}}$,
$C_0=H_{\mathrm{in}}^{-1}$, places the covariance on the curvature band
$\tfrac1\beta\Id\preceq C_0\preceq\tfrac1\alpha\Id$ (\Cref{lem:rescue-band}), and recovers the exact
optimizer in one query for Gaussian targets (\Cref{thm:gauss-recovery}). We record here the energy
quality of the rescued initializer for non-Gaussian targets, used by the global stochastic
complexity bounds.

\begin{lemma}[Rescue energy for nonquadratic targets]
\label{lem:rescue-energy}
Under \Cref{assump:logconcave-smooth}, \Cref{lem:rescue-band} places $C_0$ on the band
$[\tfrac1\beta,\tfrac1\alpha]$ pathwise, and $\Delta(a_0)$ is a deterministic finite quantity. Under
the additional Hessian-Lipschitz \Cref{assump:global-Hess-Lip}, for the point-Newton rescue
$H_{\mathrm{in}}:=\nabla^2V(m_{\mathrm{in}})$, $C_0=H_{\mathrm{in}}^{-1}$,
$m_0=m_{\mathrm{in}}-H_{\mathrm{in}}^{-1}\nabla V(m_{\mathrm{in}})$ of Phase~0
(\Cref{alg:unified}), writing $r_{\mathrm{in}}:=\norm{m_{\mathrm{in}}-m_\star}$,
\[
    \norm{m_0-m_\star}\leq\frac{\LH}{2\alpha}\Bigl(r_{\mathrm{in}}^2+\frac{N_\theta}\alpha\Bigr),
    \qquad
    \norm{C_0-C_\star}_{\Frob}\leq\frac{\sqrt{N_\theta}\,\LH}{\alpha^2}\Bigl(r_{\mathrm{in}}+\sqrt{N_\theta/\alpha}\Bigr),
\]
and hence
$\Delta(a_0)\leq\tfrac{\beta\LH^2}{8\alpha^2}(r_{\mathrm{in}}^2+\tfrac{N_\theta}\alpha)^2+\tfrac{\beta N_\theta\LH}{2\alpha^2}(r_{\mathrm{in}}+\sqrt{N_\theta/\alpha})$,
which vanishes when $\LH=0$.
\end{lemma}

\begin{proof}
\PROOFHOLE
\end{proof}
\subsubsection{Pathwise covariance bands}
\label{subsec:bands}

\begin{lemma}[Covariance bands for the rescued schemes]
\label{lem:stoch-band}
Assume \Cref{assump:logconcave-smooth} and initialize with \eqref{eq:rescue}, so
$\tfrac1\beta\Id\preceq C_0\preceq\tfrac1\alpha\Id$.
\begin{enumerate}
\item \textup{(Riemannian.)} If $0<\Delta t\leq1/\kappa$, then almost surely
$\tfrac1{2\beta}\Id\preceq C_n\preceq\tfrac1\alpha\Id$ for all $n\geq0$.
\item \textup{(KL/Bregman.)} For every $\Delta t>0$, almost surely
$\tfrac1\beta\Id\preceq C_n\preceq\tfrac1\alpha\Id$ for all $n\geq0$.
\end{enumerate}
No stopping time, clipping, or containment hypothesis is used.
\end{lemma}

\begin{proof}
\PROOFHOLE
\end{proof}

\begin{algorithm}[t]
\caption{Quadratic-rescued minibatch Price/Hessian STL Fisher--Rao scheme}
\label{alg:unified}
\begin{algorithmic}[1]
\Require initial mean $m_{\mathrm{in}}$; scheme selector $\mathsf S\in\{\Riem,\KL\}$; batch size $B$;
stepsizes $\Delta t_n>0$; iteration count $N$
\State \textbf{Phase 0 (quadratic rescue; one query).} Query the oracle at $m_{\mathrm{in}}$:
$G_{\mathrm{in}}\gets\nabla V(m_{\mathrm{in}})$, $H_{\mathrm{in}}\gets\nabla^2V(m_{\mathrm{in}})$;
set $m_0\gets m_{\mathrm{in}}-H_{\mathrm{in}}^{-1}G_{\mathrm{in}}$, $C_0\gets H_{\mathrm{in}}^{-1}$
\State \textbf{Phase I (Price/Hessian STL iterations).}
\For{$n=0,\dots,N-1$}
    \State Draw $\xi_{n,s}\overset{\text{iid}}{\sim}\N(0,\Id)$, $X_{n,s}\gets m_n+C_n^{1/2}\xi_{n,s}$, $s=1,\dots,B$; query the oracle at each $X_{n,s}$
    \State $\overline b_n\gets\frac1B\sum_s(-\nabla V(X_{n,s})+C_n^{-1/2}\xi_{n,s})$ \Comment{STL score}
    \State $\overline H_n\gets\frac1B\sum_s\nabla^2V(X_{n,s})$ \Comment{Price/Hessian covariance}
    \State $m_{n+1}\gets m_n+\Delta t_n\,C_n\overline b_n$
    \If{$\mathsf S=\Riem$}
        \State $C_{n+1}\gets C_n^{1/2}\exp\bigl(\Delta t_n(\Id-C_n^{1/2}\overline H_nC_n^{1/2})\bigr)C_n^{1/2}$
    \Else
        \State $C_{n+1}\gets(1+\Delta t_n)\bigl(C_n^{-1}+\Delta t_n\overline H_n\bigr)^{-1}$
    \EndIf
\EndFor
\end{algorithmic}
\end{algorithm}

\subsection{Intrinsic sticking-the-landing variance}
\label{sec:intrinsic-STL}

For a single sample $X\sim\rho_a$, the single-sample tangent noise is
$\xi(a;X):=(-C\varepsilon_b,\,-C\varepsilon_RC)$ with $\varepsilon_b:=\widehat b-g(a)$,
$\varepsilon_R:=-(\nabla^2V(X)-A(a))$, and the minibatch noise is
$\overline\xi_n:=\widehat\grad\calE(a_n)-\grad\calE(a_n)=\tfrac1B\sum_s\xi(a_n;X_{n,s})$.

\begin{definition}[Intrinsic Hessian fluctuation]
\label{def:Psi}
$\Psi(a):=\E_{X\sim\rho_a}\bigl[\norm{C^{1/2}(\nabla^2V(X)-A(a))C^{1/2}}_{\Frob}^2\bigr]$.
\end{definition}

\begin{lemma}[Intrinsic STL variance bound]
\label{lem:intrinsic-STL}
Under \Cref{assump:logconcave-smooth}, the STL noises are conditionally centered, and, with
$\norm{\xi(a;X)}_{a,\FR}^2=\varepsilon_b^\top C\varepsilon_b+\tfrac12\norm{C^{1/2}\varepsilon_RC^{1/2}}_{\Frob}^2$,
\begin{equation}
\label{eq:intrinsic-mean}
    \E\bigl[\varepsilon_b^\top C\varepsilon_b\bigr]\leq\norm{C^{1/2}R(a)C^{1/2}}_{\Frob}^2+\Psi(a),
    \qquad
    \E\bigl[\norm{\xi(a;X)}_{a,\FR}^2\bigr]\leq2\,\Gradnorm(a)^2+\frac32\,\Psi(a).
\end{equation}
Consequently, for every $B\geq1$,
$\E[\norm{\overline\xi_n}_{a_n,\FR}^2\mid\Fil_n]\leq B^{-1}(2\Gradnorm(a_n)^2+\tfrac32\Psi(a_n))$.
\end{lemma}

\begin{proof}
\PROOFHOLE
\end{proof}

\begin{remark}[Regularity]
\label{rem:regularity}
The differentiations above transfer to Gaussian score polynomials for observables of at most
polynomial growth, which holds under \Cref{assump:logconcave-smooth} since $\nabla V$ grows at most
linearly; no third or fourth derivatives of $V$ are used. When $\nabla^2V$ is only Lipschitz it is
differentiable a.e., and the Gaussian Poincar\'e step is applied to the Lipschitz map
$F(X)=-\nabla V(X)+C^{-1}(X-m)$ using its weak derivative $\nabla F=C^{-1}-\nabla^2V$ (equivalently by
mollification), which is all the inequality requires.
\end{remark}

\subsubsection{Assumption-minimal regime: the self-bounding estimate}
\label{sec:stoch-selfbound}

\begin{lemma}[Self-bounding estimate for the intrinsic fluctuation]
\label{lem:Psi-self-bound}
Under \Cref{assump:logconcave-smooth}, if $C\preceq\Lambda\Id$ and $\LB=\beta\Lambda$, then
\begin{equation}
\label{eq:Psi-self-bound}
    \boxed{\;
    \Psi(a)\leq2\LB N_\theta+\LB\norm{C^{1/2}R(a)C^{1/2}}_{\Frob}^2
    \leq2\LB N_\theta+2\LB\,\Gradnorm(a)^2 .
    \;}
\end{equation}
\end{lemma}

\begin{proof}
\PROOFHOLE
\end{proof}

\begin{corollary}[Self-bounded second moment]
\label{cor:second-moment}
If $C\preceq\Lambda\Id$, then for every $B\geq1$,
\begin{equation}
\label{eq:second-moment}
    \E\bigl[\norm{\widehat\grad\calE(a)}_{a,\FR}^2\mid a\bigr]
    \leq\rho_B\,\Gradnorm(a)^2+\frac{3\LB N_\theta}B,
    \qquad
    \rho_B:=1+\frac{2+3\LB}B .
\end{equation}
\end{corollary}

\begin{proof}
\PROOFHOLE
\end{proof}

\subsubsection{Hessian-Lipschitz regime}
\label{sec:stoch-Hlip}

\begin{assumption}[Global Hessian-Lipschitzness]
\label{assump:global-Hess-Lip}
There is $\LH<\infty$ with $\norm{\nabla^2V(x)-\nabla^2V(y)}_\op\leq\LH\norm{x-y}$ for all $x,y$. We
write the dimensionless parameter $\barLH:=\LH/\alpha^{3/2}$; a Gaussian target has $\LH=\barLH=0$.
\end{assumption}

\begin{lemma}[Hessian-Lipschitz control of the intrinsic fluctuation]
\label{lem:Psi-Hg}
Under \Cref{assump:global-Hess-Lip}, $\Psi(a)\leq\LH^2\norm{C}_\op(\Tr C)^2$. On the band
$C\preceq\tfrac1\alpha\Id$,
\begin{equation}
\label{eq:Psi-Hg}
    \boxed{\;\Psi(a)\leq\Psi_H:=\frac{N_\theta^2\LH^2}{\alpha^3}=N_\theta^2\barLH^2 .\;}
\end{equation}
\end{lemma}

\begin{proof}
\PROOFHOLE
\end{proof}
\subsection{Stochastic global convergence: Riemannian variant}
\label{sec:stoch-Riem}

On the rescued bands, \Cref{lem:FR-W2} gives the Fisher--Rao PL inequalities
\begin{equation}
\label{eq:PL-bands}
    \Gradnorm(a_n)^2\geq\frac1{2\kappa}\Delta_n\ \ (\text{Riemannian band }C_n\succeq\tfrac1{2\beta}\Id),
    \qquad
    \Gradnorm(a_n)^2\geq\frac1\kappa\Delta_n\ \ (\text{KL band }C_n\succeq\tfrac1\beta\Id).
\end{equation}

\begin{lemma}[Riemannian stochastic one-step descent]
\label{lem:Riem-onestep-stoch}
Assume \Cref{assump:logconcave-smooth} and $\tfrac1{2\beta}\Id\preceq C_n\preceq\tfrac1\alpha\Id$. For
one step of \eqref{eq:Riem-STL} with $0<\Delta t\leq1/\kappa$,
\begin{equation}
\label{eq:Riem-onestep-Psi}
    \E[\Delta_{n+1}\mid\Fil_n]
    \leq\Delta_n-\Delta t\Bigl(1-\kappa\Delta t\bigl(1+\tfrac2B\bigr)\Bigr)\Gradnorm(a_n)^2
    +\frac{3\kappa\Delta t^2}{2B}\Psi(a_n).
\end{equation}
\end{lemma}

\begin{proof}
\PROOFHOLE
\end{proof}

\begin{theorem}[Riemannian, Hessian-Lipschitz]
\label{thm:Riem-Hlip}
Assume \Cref{assump:logconcave-smooth,assump:global-Hess-Lip}, initialize with \eqref{eq:rescue}, and
run \eqref{eq:Riem-STL} with $B\geq1$ and $\Delta t=\tfrac1{2\kappa(1+2/B)}$. Then for every $N\geq0$,
\begin{equation}
\label{eq:Riem-Hlip}
    \boxed{\;
    \E\Delta_N\leq\Bigl(1-\frac1{8\kappa^2(1+2/B)}\Bigr)^N\Delta(a_0)+\frac{3\kappa}{B+2}\,N_\theta^2\barLH^2 .
    \;}
\end{equation}
For $B=1$: $\Delta t=\tfrac1{6\kappa}$, $\E\Delta_N\leq(1-\tfrac1{24\kappa^2})^N\Delta(a_0)+\kappa N_\theta^2\barLH^2$,
so the single-sample iteration and oracle complexities are $N_{\Riem}=\widetilde O(\kappa^2)$ and
$Q_{\Riem}=1+\widetilde O(\kappa^2)$.
\end{theorem}

\begin{proof}
\PROOFHOLE
\end{proof}

\begin{theorem}[Riemannian, assumption-minimal self-bounding]
\label{thm:Riem-selfbound}
Assume \Cref{assump:logconcave-smooth}, initialize with \eqref{eq:rescue}, and run \eqref{eq:Riem-STL}
with $B\geq1$ and $\Delta t=\tfrac1{2\kappa\rho_B}$, $\rho_B:=1+\tfrac{2+3\kappa}B$. Then for every
$N\geq0$,
\begin{equation}
\label{eq:Riem-selfbound}
    \boxed{\;
    \E\Delta_N\leq\Bigl(1-\frac1{8\kappa^2\rho_B}\Bigr)^N\Delta(a_0)+\frac{6\kappa^2N_\theta}{B+2+3\kappa} .
    \;}
\end{equation}
The floor is $O(\kappa N_\theta)$ for every $B\lesssim\kappa$ and at most $2\kappa N_\theta$ for every
$B\geq1$. For $B=1$, $\rho_1=3+3\kappa\asymp\kappa$, giving $N_{\Riem}=\widetilde O(\kappa^3)$; for
$B=\Theta(\kappa)$, $\rho_B=O(1)$, giving $N_{\Riem}=\widetilde O(\kappa^2)$ rounds and
$Q_{\Riem}=BN_{\Riem}=\widetilde O(\kappa^3)$ oracle pairs. No Hessian-Lipschitz assumption is used.
\end{theorem}

\begin{proof}
\PROOFHOLE
\end{proof}
\subsection{Stochastic global convergence: KL/Bregman variant}
\label{sec:stoch-KL}

The KL/Bregman scheme \eqref{eq:KL-STL} enjoys the tighter pathwise band of
\Cref{lem:stoch-band}(2), $\tfrac1\beta\Id\preceq C_n\preceq\tfrac1\alpha\Id$ for \emph{every}
$\Delta t>0$, and the corresponding Fisher--Rao PL inequality from \eqref{eq:PL-bands},
$\Gradnorm(a_n)^2\geq\tfrac1\kappa\Delta_n$. The key observation is that the implicit resolvent
covariance update is exactly an affine-invariant exponential retraction, so the one-step estimate
follows from the arbitrary-direction retraction smoothness of \Cref{lem:retr-smooth} without any
principal-square-root coupling.

\begin{lemma}[Logarithmic tangent of the KL resolvent update]
\label{lem:KL-logtangent}
Fix $a=(m,C)$ with $\tfrac1\beta\Id\preceq C\preceq\tfrac1\alpha\Id$ and $\alpha\Id\preceq\overline H\preceq\beta\Id$,
and set $P:=C^{1/2}\overline HC^{1/2}$, so that $\tfrac1\kappa\Id\preceq P\preceq\kappa\Id$. Write
$\overline S:=\Id-P$ and define the logarithmic tangent
\begin{equation}
\label{eq:Wlog}
    W(P):=\log(1+\Delta t)\,\Id-\log(\Id+\Delta t\,P).
\end{equation}
Then the KL covariance update $C^+=(1+\Delta t)(C^{-1}+\Delta t\,\overline H)^{-1}$ satisfies
\begin{equation}
\label{eq:KL-retr}
    C^+=C^{1/2}e^{W(P)}C^{1/2},
    \qquad\text{equivalently}\qquad
    a^+=\Retr_a\bigl(\Delta t\,C\overline b,\;C^{1/2}W(P)C^{1/2}\bigr),
\end{equation}
where $\overline b$ is the minibatch STL score. Moreover, with $\mathcal R(P):=W(P)-\Delta t\,\overline S$,
\begin{equation}
\label{eq:Wlog-bounds}
    \norm{W(P)}_\op\leq\Delta t\,\kappa,
    \qquad
    \norm{W(P)}_{\Frob}\leq\Delta t\,\norm{\overline S}_{\Frob},
    \qquad
    \norm{\mathcal R(P)}_{\Frob}\leq\kappa\,\Delta t^2\,\norm{\overline S}_{\Frob}.
\end{equation}
In particular, for $\Delta t\leq1/(8\kappa)$, $\norm{W(P)}_\op\leq\tfrac18<1$, so the tangent
$\zeta:=(\Delta t\,C\overline b,\,C^{1/2}W(P)C^{1/2})$ satisfies the scope condition \eqref{eq:scope}.
\end{lemma}

\begin{proof}
\PROOFHOLE
\end{proof}

\begin{lemma}[KL stochastic one-step descent]
\label{lem:KL-onestep-stoch}
Assume \Cref{assump:logconcave-smooth} and $\tfrac1\beta\Id\preceq C_n\preceq\tfrac1\alpha\Id$. For
one step of \eqref{eq:KL-STL} with $0<\Delta t\leq1/(8\kappa)$,
\begin{equation}
\label{eq:KL-onestep-Psi}
    \E[\Delta_{n+1}\mid\Fil_n]
    \leq\Delta_n-\frac{\Delta t}2\,\Gradnorm(a_n)^2
    +\frac{7\kappa\,\Delta t^2}{4B}\,\Psi(a_n).
\end{equation}
\end{lemma}

\begin{proof}
\PROOFHOLE
\end{proof}

\begin{theorem}[KL, Hessian-Lipschitz]
\label{thm:KL-Hlip}
Assume \Cref{assump:logconcave-smooth,assump:global-Hess-Lip}, initialize with \eqref{eq:rescue},
and run \eqref{eq:KL-STL} with $B\geq1$ and $\Delta t=\tfrac1{8\kappa}$. Then for every $N\geq0$,
\begin{equation}
\label{eq:KL-Hlip}
    \boxed{\;
    \E\Delta_N\leq\Bigl(1-\frac1{16\kappa^2}\Bigr)^N\Delta(a_0)+\frac{7\kappa}{16B}\,N_\theta^2\barLH^2 .
    \;}
\end{equation}
For $B=1$ the single-sample iteration and oracle complexities are $N_{\KL}=\widetilde O(\kappa^2)$ and
$Q_{\KL}=1+\widetilde O(\kappa^2)$.
\end{theorem}

\begin{proof}
\PROOFHOLE
\end{proof}

\begin{theorem}[KL, assumption-minimal self-bounding]
\label{thm:KL-selfbound}
Assume \Cref{assump:logconcave-smooth}, initialize with \eqref{eq:rescue}, and run \eqref{eq:KL-STL}
with $B\geq1$ and $\Delta t=\min\{\tfrac1{8\kappa},\tfrac B{14\kappa^2}\}$. Then for every $N\geq0$,
\begin{equation}
\label{eq:KL-selfbound}
    \boxed{\;
    \E\Delta_N\leq\Bigl(1-\frac{\Delta t}{4\kappa}\Bigr)^N\Delta(a_0)+14\kappa^3\Delta t\,\frac{N_\theta}B .
    \;}
\end{equation}
For $B\leq\tfrac74\kappa$, $\Delta t=\tfrac B{14\kappa^2}$, the floor is exactly $\kappa N_\theta$ and
the contraction is $1-\tfrac B{56\kappa^3}$; for $B=1$ this gives $N_{\KL}=\widetilde O(\kappa^3)$, and
for $B=\Theta(\kappa)$, $N_{\KL}=\widetilde O(\kappa^2)$ rounds and $Q_{\KL}=\widetilde O(\kappa^3)$
oracle pairs. No Hessian-Lipschitz assumption is used.
\end{theorem}

\begin{proof}
\PROOFHOLE
\end{proof}

\begin{remark}[The two schemes side by side]
\label{rem:KL-vs-Riem}
On the rescued bands the two schemes attain the same leading complexity in each regime:
$\widetilde O(\kappa^2)$ single-sample iterations under \Cref{assump:global-Hess-Lip} and
$\widetilde O(\kappa^3)$ under \Cref{assump:logconcave-smooth} alone, both unconditionally. The KL band
$\tfrac1\beta\Id\preceq C_n\preceq\tfrac1\alpha\Id$ holds for every $\Delta t>0$ and gives the
stronger PL constant $\tfrac1\kappa$ (versus $\tfrac1{2\kappa}$ for the Riemannian band
$\tfrac1{2\beta}\Id\preceq C_n$), while the Riemannian scheme requires $\Delta t\leq1/\kappa$ to
maintain its band. The structural reason the KL analysis matches the Riemannian one is
\Cref{lem:KL-logtangent}: the implicit resolvent update is the exponential retraction
$C^{1/2}e^{W(P)}C^{1/2}$ with $W(P)$ a spectral function of the sampled relative Hessian $P$, so its
first-order stochastic part is unbiased and only a second-order fluctuation survives. Both floors are
proportional to $\Psi_H=N_\theta^2\barLH^2$ (Hessian-Lipschitz) or $\kappa N_\theta$ (self-bounding),
and both vanish at the estimator level for Gaussian targets (\Cref{prop:gauss-nofloor}).
\end{remark}
\subsection{Stochastic local convergence: Gaussian-core Price/Hessian--STL theory}
\label{sec:stoch-local}

The global stochastic theorems (\Cref{sec:stoch-Riem,sec:stoch-KL}) drive $\E\Delta_N$ to a
$\kappa$-dependent floor at the global $\kappa$-dependent rate. Once an iterate has entered a
certified local region $\Ureg_{\delta_\bullet}$, however, the \emph{deterministic} discrete local
map contracts at the rate of \Cref{thm:disc-local}. The local linearized spectral scale is
$4+\Gamma=O(1+N_\theta\log\kappa)$ (after whitening, independent of the global covariance
conditioning); together with the universal certified stepsize $\Delta t=\Theta((4+\Gamma)^{-1})$ in the
large-batch branch $B\gtrsim V_1$ ($V_1=24+6N_\theta\LHs^2$) this yields the
$O((4+\Gamma)^2\log\tfrac1\epsilon)$ iteration bound derived below (\Cref{rem:local-condnum}),
quadratic in the scale (for general $B$, with the extra factor $\max\{1,V_1/B\}$). This section gives the stochastic analogue of \Cref{thm:disc-local}: a mean-square
contraction on the non-exit event, i.e.\ a bound on $\E[\norm{a_N-a_\star}_{\star,\bullet}^2;\tau>N]$,
with a \emph{target-adaptive} additive floor proportional to the \emph{squared} whitened
Hessian-Lipschitz modulus $\LHs^2$, together with a finite-horizon high-probability containment bound. This is a
\emph{killed-process} quantity (trajectories that have exited by time $N$ contribute zero), not the
mean square of the stopped process $a_{N\wedge\tau}$; the containment proposition supplies the exit
control that makes it operative. As already noted in \Cref{prop:entry-residual}, a gate alone yields
neither pathwise nor expected containment for the stochastic schemes; the results here are the
promised separate stability argument. The standing assumptions do not guarantee bounded support for
the STL mean-score noise: even for a Gaussian target, away from $C=\Id$ it equals
$(C^{-1/2}-C^{1/2})Z$ (see the Gaussian example below), which is unbounded---this single counterexample
already shows that pathwise invariance of a bounded local energy sublevel cannot be guaranteed under
the standing assumptions at fixed stepsize without projection, clipping, or safeguarding, so the
honest statements are non-exit contraction plus exit-probability control rather than infinite-horizon
containment. For Gaussian targets $\LHs=0$, the
\emph{additive} floor vanishes (\Cref{cor:stoch-local-gauss}); the noise is pathwise zero (deterministic
contraction, no exit control) whenever $C_0=C_\star$, even if $m_0\neq m_\star$, with quadratic rescue
the special case $a_0=a_\star$, while generic $C_0\neq C_\star$ retains unbounded mean noise.

\subsubsection{Whitened reduction and the local modulus}
\label{subsec:local-whiten}

The Fisher--Rao and KL natural-gradient schemes are equivariant under invertible affine
reparametrization of $\R^{N_\theta}$. Whitening by $C_\star^{1/2}$ maps $a_\star$ to $(0,\Id)$ and
$V$ to the whitened potential $\widehat V(z)=V(m_\star+C_\star^{1/2}z)$ of \Cref{sec:local}, under
which $\alpha_\star\Id\preceq\nabla^2\widehat V\preceq\beta_\star\Id$, $A(a_\star)=\Id$, and the
minibatch updates \eqref{eq:Riem-STL}--\eqref{eq:KL-STL} take the same form with $V\to\widehat V$,
$\nabla^2 V\to\nabla^2\widehat V$. We work in these coordinates throughout the section; $C$ denotes
the (whitened) covariance and $\overline H_n=\tfrac1B\sum_s\nabla^2\widehat V(X_{n,s})$, so that
$\alpha_\star\Id\preceq\overline H_n\preceq\beta_\star\Id$ almost surely by \eqref{eq:pathwise-Hess}.
Write
\[
    \overline\varepsilon_{b,n}:=\overline b_n-g(a_n),\qquad
    \overline\varepsilon_{H,n}:=\overline H_n-A(a_n),
\]
which are conditionally centered by \eqref{eq:STL-unbiased}; at the single-sample level
$\varepsilon_b:=\bigl(-\nabla\widehat V(X)+C^{-1}(X-m)\bigr)-g(a)$ and
$\varepsilon_H:=\nabla^2\widehat V(X)-A(a)=-\varepsilon_R$, so that
$\E\norm{\varepsilon_H}_{\Frob}^2$ is the (un-whitened-weighted) form of the intrinsic fluctuation
$\Psi$ of \Cref{def:Psi}.

\begin{lemma}[Whitened local Hessian-Lipschitz modulus]
\label{lem:whitened-HessLip}
Under \Cref{assump:global-Hess-Lip} (global modulus $\LH$), the whitened potential $\widehat V$ has a
Hessian-Lipschitz modulus $\LHs$, i.e.\
$\norm{\nabla^2\widehat V(x)-\nabla^2\widehat V(y)}_{\op}\leq\LHs\norm{x-y}$ for all $x,y$, with
\[
    \LHs\leq\LH\,\norm{C_\star}_{\op}^{3/2}\leq\LH\,\alpha^{-3/2}=\barLH .
\]
\end{lemma}

\begin{proof}
\PROOFHOLE
\end{proof}

\begin{definition}[Stochastically admissible level]
\label{def:stoch-admissible}
For a fixed scheme $\bullet\in\{\Riem,\KL\}$ and a fixed stepsize $\Delta t$, a level $\delta_\bullet>0$
is \emph{stochastically admissible} (for that $\Delta t$) if the discrete radius condition
$\Kdisc(\delta_\bullet)\,\mathfrak b_{\star,\bullet}(\delta_\bullet)\sqrt{\delta_\bullet}\leq\gamma_\bullet$
of \Cref{thm:disc-local} holds and additionally
\[
    \ell_{\delta_\bullet}\geq\tfrac12,\qquad U_{\delta_\bullet}\leq\tfrac32,\qquad
    \mathfrak b_{\star,\bullet}(\delta_\bullet)\sqrt{\delta_\bullet}\leq1 .
\]
This is a joint condition on the pair $(\Delta t,\delta_\bullet)$: $\Kdisc$, $\gamma_\bullet$ (for KL),
and $\mathfrak b_{\star,\bullet}$ all depend on $\Delta t$, so it is understood that $\Delta t$ is fixed
first---satisfying the linearized-gap and stochastic-absorption bounds of
\Cref{thm:stoch-local-Riem,thm:stoch-local-KL}---and $\delta_\bullet$ chosen afterward.
\end{definition}

\begin{remark}[Existence]
\label{rem:stoch-admissible-exists}
As $\delta\downarrow0$, the band roots satisfy $\ell_\delta\uparrow1$ and $U_\delta\downarrow1$
(\Cref{cor:band}), $\mathfrak b_{\star,\bullet}(\delta)\sqrt\delta\to0$, and, \emph{for every fixed
admissible stepsize $\Delta t>0$}, $\Kdisc(\delta)\to\Kdisc(0)<\infty$ as $\delta\downarrow0$ (the
supremum of the continuous $\Delta t^{-1}D^2T_{\bullet,\Delta t}$ over the shrinking hull
$\Hreg_\delta\to\{a_\star\}$, finite by \Cref{rem:Kdisc-finite}), so the discrete radius condition
$\Kdisc(\delta)\mathfrak b_{\star,\bullet}(\delta)\sqrt\delta\leq\gamma_\bullet$ holds for all small
$\delta$. Hence, for every fixed admissible $\Delta t$, sufficiently small stochastically admissible
levels $\delta_\bullet$ exist. On $\Ureg_{\delta_\bullet}$ one has
$\tfrac12\Id\preceq C\preceq\tfrac32\Id$ by \Cref{cor:band}(i), which is the band hypothesis of the
variance lemma below. This is an \emph{existence} statement: the manuscript does not give a
closed-form value of the switching level $\delta_\bullet$ in terms of $\alpha,\beta,N_\theta,\LH$
alone, since $\Kdisc$ is defined through a supremum of $D^2T_{\bullet,\Delta t}$ over the local hull
and is not bounded in closed form here; an explicit upper bound on $\Kdisc$ (a natural later
strengthening) would make the discrete switching threshold precomputable.
\end{remark}

\subsubsection{Local Price/Hessian--STL variance}
\label{subsec:local-var}

\begin{lemma}[Local variance decomposition]
\label{lem:local-variance}
Assume \Cref{assump:logconcave-smooth,assump:global-Hess-Lip}, work in whitened coordinates, and let
$a=(m,C)$ with $\tfrac12\Id\preceq C\preceq\tfrac32\Id$. Then for a single sample $X\sim\rho_a$,
\begin{equation}
\label{eq:local-var}
    \E\Bigl[\norm{\varepsilon_b}^2+\tfrac12\norm{\varepsilon_H}_{\Frob}^2\Bigr]
    \leq V_0+V_1\,\norm{a-a_\star}_\star^2,
    \qquad
    V_0:=3N_\theta^2\LHs^2,\quad V_1:=24+6N_\theta\LHs^2 .
\end{equation}
By conditional independence, the minibatch noise obeys
$\E[\norm{\overline\varepsilon_b}^2+\tfrac12\norm{\overline\varepsilon_H}_{\Frob}^2\mid\Fil]\leq B^{-1}(V_0+V_1\norm{a-a_\star}_\star^2)$.
At the optimizer, $\E[\norm{\varepsilon_b(a_\star)}^2+\tfrac12\norm{\varepsilon_H(a_\star)}_{\Frob}^2]\leq\sigma_\star^2:=\tfrac32N_\theta^2\LHs^2$,
which is $0$ for Gaussian targets.
\end{lemma}

\begin{proof}
\PROOFHOLE
\end{proof}

\subsubsection{One-step stochastic perturbation of the local maps}
\label{subsec:local-perturb}

We compare one stochastic step to the deterministic local step $T_\bullet$ of \Cref{thm:disc-local}.
The following elementary bounds on the matrix exponential are the only nonlinear ingredient.

\begin{lemma}[Exponential perturbation bounds]
\label{lem:exp-perturb}
Let $S,\widehat S\in\Sym(N_\theta)$ with $S\preceq\Id$, $\widehat S\preceq\Id$, put
$\Delta:=\widehat S-S$, and let $0<\Delta t\leq1$. Then
\begin{align*}
\textup{(i)}\quad & \norm{e^{\Delta t\widehat S}-e^{\Delta tS}}_{\Frob}\leq e\,\Delta t\,\norm{\Delta}_{\Frob},\\
\textup{(ii)}\quad & \norm{e^{\Delta t\widehat S}-e^{\Delta tS}-D\exp(\Delta tS)[\Delta t\Delta]}_{\Frob}
\leq\tfrac{e}{2}\,\Delta t^2\,\norm{\Delta}_{\op}\norm{\Delta}_{\Frob}.
\end{align*}
\end{lemma}

\begin{proof}
\PROOFHOLE
\end{proof}

\begin{lemma}[One-step perturbation, Riemannian]
\label{lem:local-perturb-Riem}
Assume \Cref{assump:logconcave-smooth,assump:global-Hess-Lip}, work in whitened coordinates, and let
$a=(m,C)$ with $\tfrac12\Id\preceq C\preceq\tfrac32\Id$ and $\norm{T_{\Riem}(a)-a_\star}_\star\leq1$,
where $T_{\Riem}$ is the deterministic map of \eqref{eq:Riem-det} and $\widehat a^+$ the one stochastic
step \eqref{eq:Riem-STL} with batch $B$ and $0<\Delta t\leq1$. Then
\[
    \E\bigl[\norm{\widehat a^+-a_\star}_\star^2\mid\Fil\bigr]
    \leq\norm{T_{\Riem}(a)-a_\star}_\star^2+\frac{64\,\Delta t^2}{B}\bigl(V_0+V_1\norm{a-a_\star}_\star^2\bigr).
\]
\end{lemma}

\begin{proof}
\PROOFHOLE
\end{proof}

\begin{lemma}[One-step perturbation, KL]
\label{lem:local-perturb-KL}
Under the hypotheses of \Cref{lem:local-perturb-Riem} but for the KL maps
\eqref{eq:KL-det}/\eqref{eq:KL-STL} and the norm $\norm\cdot_{\star,\Delta t}$ (with
$\norm{T_{\KL}(a)-a_\star}_{\star,\Delta t}\leq1$),
\[
    \E\bigl[\norm{\widehat a^+-a_\star}_{\star,\Delta t}^2\mid\Fil\bigr]
    \leq\norm{T_{\KL}(a)-a_\star}_{\star,\Delta t}^2+\frac{68\,\Delta t^2}{B}\bigl(V_0+V_1\norm{a-a_\star}_\star^2\bigr).
\]
\end{lemma}

\begin{proof}
\PROOFHOLE
\end{proof}

\subsubsection{Localized non-exit mean-square contraction}
\label{subsec:local-contraction}

\begin{theorem}[Stochastic local convergence, Riemannian]
\label{thm:stoch-local-Riem}
Assume \Cref{assump:logconcave-smooth,assump:global-Hess-Lip}, work in whitened coordinates, and set
$A_{\Riem}:=4+\Gamma$, $\mu_{\Riem}:=1/(2A_{\Riem})$, $V_1=24+6N_\theta\LHs^2$. Fix a stepsize
\[
    0<\Delta t\leq\min\Bigl\{\tfrac{2}{4+\Gamma},\ \tfrac{B}{256\,A_{\Riem}\,V_1}\Bigr\},
\]
and let $\delta_{\Riem}$ be stochastically admissible for this $\Delta t$ (\Cref{def:stoch-admissible}).
Run \eqref{eq:Riem-STL} with batch $B$.
Let $\tau_{\Riem}:=\inf\{n\geq0:a_n\notin\Ureg_{\delta_{\Riem}}\}$ and $a_0\in\Ureg_{\delta_{\Riem}}$.
Then for all $N\geq0$,
\[
    \boxed{\;
    \E\bigl[\norm{a_N-a_\star}_\star^2\,\mathbf 1_{\{\tau_{\Riem}>N\}}\bigr]
    \leq\bigl(1-\mu_{\Riem}\Delta t\bigr)^N\norm{a_0-a_\star}_\star^2
    +\frac{384\,\Delta t\,(4+\Gamma)\,N_\theta^2\LHs^2}{B}.
    \;}
\]
The left-hand side is the \emph{non-exit} (killed) second moment: trajectories that leave
$\Ureg_{\delta_{\Riem}}$ by time $N$ contribute zero, so this is not the mean square of the stopped
process $a_{N\wedge\tau_{\Riem}}$; \Cref{prop:finite-horizon} bounds the discarded mass
$\Prob(\tau_{\Riem}\leq N)$, and the two together give an operative guarantee.
\end{theorem}

\begin{proof}
\PROOFHOLE
\end{proof}

\begin{theorem}[Stochastic local convergence, KL]
\label{thm:stoch-local-KL}
Under the analogous hypotheses, set $A_{\KL}:=4+2\Delta t+\Gamma$, $\mu_{\KL}:=1/(2A_{\KL})$. Fix a
stepsize
\[
    0<\Delta t\leq\min\Bigl\{\tfrac2{4+\Gamma},\ \tfrac{B}{272\,A_{\KL}\,V_1}\Bigr\},
\]
let $\delta_{\KL}$ be stochastically admissible for this $\Delta t$, and put $\tau_{\KL}:=\inf\{n:a_n\notin\Ureg_{\delta_{\KL}}\}$;
running \eqref{eq:KL-STL} with batch $B$, one has for $a_0\in\Ureg_{\delta_{\KL}}$ and all $N\geq0$,
\[
    \boxed{\;
    \E\bigl[\norm{a_N-a_\star}_{\star,\Delta t}^2\,\mathbf 1_{\{\tau_{\KL}>N\}}\bigr]
    \leq\bigl(1-\mu_{\KL}\Delta t\bigr)^N\norm{a_0-a_\star}_{\star,\Delta t}^2
    +\frac{408\,\Delta t\,(4+2\Delta t+\Gamma)\,N_\theta^2\LHs^2}{B}.
    \;}
\]
As in the Riemannian case the left-hand side is the non-exit (killed) second moment, paired with the
exit bound of \Cref{prop:finite-horizon}.
\end{theorem}

\begin{proof}
\PROOFHOLE
\end{proof}

\subsubsection{Finite-horizon containment}
\label{subsec:finite-horizon}

The non-exit bounds contract the surviving mass on $\{\tau_\bullet>N\}$; to control the discarded mass
we bound the exit probability. Fix a radius $r_\bullet>0$ with $r_\bullet\leq\tfrac1{2\sqrt2}$ and
$\{a:\norm{a-a_\star}_{\star,\bullet}\leq r_\bullet\}\subseteq\Ureg_{\delta_\bullet}$ (both hold for
$r_\bullet$ small, the first forcing $\tfrac12\Id\preceq C\preceq\tfrac32\Id$ since then
$\norm{C-\Id}_{\Frob}\leq\sqrt2\,r_\bullet\leq\tfrac12$), and let
$\tau_{r_\bullet}:=\inf\{n\geq0:\norm{a_n-a_\star}_{\star,\bullet}>r_\bullet\}$.

\begin{proposition}[Exit probability]
\label{prop:finite-horizon}
In the setting of \Cref{thm:stoch-local-Riem} (resp.\ \Cref{thm:stoch-local-KL}), with
$q_\bullet:=1-\mu_\bullet\Delta t$ and $\eta_\bullet$ the corresponding one-step floor increment
($\eta_{\Riem}=192\Delta t^2N_\theta^2\LHs^2/B$, $\eta_{\KL}=204\Delta t^2N_\theta^2\LHs^2/B$), if
$a_0$ satisfies $\norm{a_0-a_\star}_{\star,\bullet}\leq r_\bullet$ then for every $N\geq1$,
\[
    \Prob\bigl(\tau_{r_\bullet}\leq N\bigr)
    \leq\frac{q_\bullet\,\norm{a_0-a_\star}_{\star,\bullet}^2+N\,\eta_\bullet}{(1-q_\bullet)\,r_\bullet^2}
    =\frac{q_\bullet\,\norm{a_0-a_\star}_{\star,\bullet}^2+N\,\eta_\bullet}{\mu_\bullet\,\Delta t\,r_\bullet^2}.
\]
The two terms are controlled separately: the batch $B$ affects only $\eta_\bullet$, not the
initial-condition term $q_\bullet\norm{a_0-a_\star}_{\star,\bullet}^2/(\mu_\bullet\Delta t\,r_\bullet^2)$,
so containment needs both a \emph{deep-entry} condition on $a_0$ and a batch condition. Indeed
$\Prob(\tau_{r_\bullet}\leq N)\leq\epsilon'$ holds whenever
\[
    \norm{a_0-a_\star}_{\star,\bullet}^2\leq\frac{\epsilon'}2\,\frac{\mu_\bullet\Delta t\,r_\bullet^2}{q_\bullet}
    \qquad\text{and}\qquad
    N\eta_\bullet\leq\frac{\epsilon'}2\,\mu_\bullet\Delta t\,r_\bullet^2,
\]
the latter, on substituting $\eta_\bullet=c'_\bullet\Delta t^2N_\theta^2\LHs^2/B$ ($c'_{\Riem}=192$,
$c'_{\KL}=204$), becoming
$\dfrac{c'_\bullet N\Delta t^2N_\theta^2\LHs^2}{B}\leq\dfrac{\epsilon'}2\mu_\bullet\Delta t\,r_\bullet^2$,
i.e.\
\[
    B\geq\frac{2c'_\bullet\,N\,\Delta t\,N_\theta^2\LHs^2}{\epsilon'\,\mu_\bullet\,r_\bullet^2}
    =\frac{c_\bullet\,N\,\Delta t\,N_\theta^2\LHs^2}{\epsilon'\,\mu_\bullet\,r_\bullet^2},
    \qquad c_{\Riem}=384,\quad c_{\KL}=408.
\]
The deep-entry requirement motivates two nested regions
$\mathcal U_{\mathrm{in}}\subsetneq\mathcal U_{\mathrm{out}}$: the local phase is entered on the smaller
$\mathcal U_{\mathrm{in}}$ (guaranteeing small $\norm{a_0-a_\star}_{\star,\bullet}$) while the exit time
is measured against the larger $\mathcal U_{\mathrm{out}}$ (radius $r_\bullet$); \Cref{cor:stoch-three-stage}
supplies the high-probability entry into $\mathcal U_{\mathrm{in}}$.
\end{proposition}

\begin{proof}
\PROOFHOLE
\end{proof}

\subsubsection{Local rate, condition number, and the Gaussian limit}
\label{subsec:local-condnum}

\begin{remark}[Condition-number dependence and two stepsize regimes]
\label{rem:local-condnum}
By $\kappa_\star\leq\kappa^2$ (so $\log\kappa_\star\leq2\log\kappa$) and the definition of $\Gamma$,
\[
    \Gamma\leq N_\theta\Bigl(3+\tfrac4{\sqrt\pi}(1+2\log\kappa)\Bigr)=O(1+N_\theta\log\kappa),
    \qquad A_\bullet=4+O(\Delta t)+\Gamma=O(1+N_\theta\log\kappa).
\]
Crucially, \Cref{lem:local-map-contraction} removed the energy-descent bound
$\Delta t\leq1/(\beta_\star\bar\Lambda_\delta)$ from the stochastic local stepsize condition: since
$\bar\Lambda_\delta\geq1/\alpha_\star$ gives $\beta_\star\bar\Lambda_\delta\geq\kappa_\star$ (up to
$\kappa^2$), retaining that bound would have forced $\Delta t=O(\kappa^{-2})$ and contaminated the
local rate with a spurious polynomial-$\kappa$ factor. Write $A_0:=4+\Gamma$, so $A_{\Riem}=A_0$ and,
since $\Delta t\leq2/A_0\leq\tfrac12$ (as $A_0\geq4$), $A_{\KL}=A_0+2\Delta t\leq A_0+1$. The corrected
windows are
\[
    0<\Delta t\leq\min\Bigl\{\tfrac2{A_0},\ \tfrac{B}{256\,A_0\,V_1}\Bigr\}\ (\Riem),
    \qquad
    0<\Delta t\leq\min\Bigl\{\tfrac2{A_0},\ \tfrac{B}{272\,A_{\KL}\,V_1}\Bigr\}\ (\KL),
\]
with $V_1=24+6N_\theta\LHs^2$; the first (linearized-gap) restriction is $\Delta t\leq2/(4+\Gamma)$ for
both schemes. The Riemannian choice below is the largest value permitted by the two displayed
sufficient conditions; the KL choice is an explicit \emph{conservative} certified stepsize that avoids
the implicit dependence of $A_{\KL}$ on $\Delta t$ by replacing $A_{\KL}$ with $A_0+1$:
\[
    \Delta t_{\Riem}^\star:=\min\Bigl\{\tfrac2{A_0},\ \tfrac{B}{256\,A_0\,V_1}\Bigr\},
    \qquad
    \Delta t_{\KL}^\star:=\min\Bigl\{\tfrac2{A_0},\ \tfrac{B}{272\,(A_0+1)\,V_1}\Bigr\}\leq\tfrac{B}{272\,A_{\KL}\,V_1}.
\]
(The exact largest KL stepsize solves $2\Delta t^2+A_0\Delta t-\tfrac{B}{272V_1}\leq0$, i.e.\
$\Delta t\leq\min\{\tfrac2{A_0},\,\tfrac14(\sqrt{A_0^2+B/(34V_1)}-A_0)\}$; the conservative choice above is
simpler and of the same order.)
At $\Delta t_\bullet^\star$ the per-step non-exit factor is
$q_\bullet=1-\mu_\bullet\Delta t_\bullet^\star=1-\Theta\!\bigl(A_0^{-2}\min\{1,B/V_1\}\bigr)$ for both
schemes (using $\mu_{\KL}\Delta t_{\KL}^\star\geq\Delta t_{\KL}^\star/(2(A_0+1))$). The floor
$F_\bullet=c'_\bullet\Delta t\,A_\bullet N_\theta^2\LHs^2/B$ ($c'_{\Riem}=384$, $c'_{\KL}=408$)
evaluates, via $\Delta t_\bullet^\star A_\bullet=\Theta(\min\{2,B/(c_\bullet V_1)\})$ ($c_{\Riem}=256$,
$c_{\KL}=272$), to
\[
    F_\bullet=O\!\Bigl(N_\theta^2\LHs^2\min\{\tfrac1B,\tfrac1{V_1}\}\Bigr)
    =\begin{cases}O(N_\theta^2\LHs^2/B), & B\gtrsim V_1,\\[1mm] O(N_\theta^2\LHs^2/V_1), & B\lesssim V_1,\end{cases}
\]
so the simplified floor $O(N_\theta^2\LHs^2/B)$ holds precisely in the large-batch branch. Hence, for
both schemes,
\[
    N_{\mathrm{loc}}=\frac1{1-q_\bullet}\log\frac{r_0^2}\epsilon
    =O\!\Bigl((4+\Gamma)^2\max\{1,V_1/B\}\log\tfrac{r_0^2}\epsilon\Bigr).
\]
Two regimes result.
\emph{(i) Small-$\Gamma$ regime} (which includes sufficiently near-Gaussian targets). If $\Gamma=O(1)$
(e.g.\ $\beta_\star-1=O(1)$) and $V_1/B=O(1)$ (batch $B\gtrsim V_1$), then $A_\bullet=O(1)$,
$\Delta t_\bullet^\star=\Theta(1)$, and $N_{\mathrm{loc}}=O(\log\tfrac1\epsilon)$, \emph{independent of
$\kappa$}. (This no longer requires
$\beta_\star\bar\Lambda_\delta=O(1)$, which the removed constraint would have demanded; $\Gamma=O(1)$
alone does not bound $V_1$ or the admissible radius, so the batch condition $B\gtrsim V_1$ is stated
separately.)
\emph{(ii) Universal window.} When $\Gamma$ grows, the worst case $\Gamma=\Theta(N_\theta\log\kappa)$
gives, in the large-batch branch $B\gtrsim V_1$,
$N_{\mathrm{loc}}=O\bigl((4+\Gamma)^2\log\tfrac1\epsilon\bigr)=O\bigl((N_\theta\log\kappa)^2\log\tfrac1\epsilon\bigr)$;
for general $B$ this carries the extra factor $\max\{1,V_1/B\}$.
Thus, suppressing dimension and the non-Gaussianity factor $V_1/B$, the universal local complexity is
$O((\log\kappa)^2\log\tfrac1\epsilon)$. A genuinely intermediate $O(N_\theta\log\kappa\cdot\log\tfrac1\epsilon)$
bound, valid when $\Gamma=\Theta(\log\kappa)$ while $\Delta t=\Theta(1)$, is \emph{not} proved here: it
would require a sharper $\Gamma$-independent linearized-gap stepsize theorem. This is the same
limitation flagged in the deterministic three-stage analysis and is separate from the
polynomial-$\kappa$ global floor.
\end{remark}

\begin{corollary}[Gaussian targets: zero additive floor; deterministic under matched covariance; exact recovery under rescue]
\label{cor:stoch-local-gauss}
If $V(x)=\tfrac12(x-\mu)^\top H(x-\mu)+c$ with $H\succ0$, then $\LHs=0$, so $V_0=\sigma_\star^2=0$ and
$\eta_\bullet=0$. Under the stepsize hypotheses of the corresponding stochastic-local theorem
(\Cref{thm:stoch-local-Riem,thm:stoch-local-KL})---in particular $0<\Delta t\leq2/(4+\Gamma)=\tfrac12$
for Gaussian targets, since $\Gamma=0$, which gives $|1-\Delta t|<1$---three cases are distinguished by
the covariance initialization (whitened, so $C_\star=\Id$).
\begin{enumerate}
\item[(a)] \emph{Generic $C_0\neq\Id$.} The \emph{additive} floors in
\Cref{thm:stoch-local-Riem,thm:stoch-local-KL} vanish, and the exit bound of
\Cref{prop:finite-horizon} loses its $N$-growing term:
\[
    \E[\norm{a_N-a_\star}_{\star,\bullet}^2;\tau_\bullet>N]\leq(1-\mu_\bullet\Delta t)^N\norm{a_0-a_\star}_{\star,\bullet}^2,
    \qquad
    \Prob(\tau_{r_\bullet}\leq N)\leq\frac{q_\bullet\norm{a_0-a_\star}_{\star,\bullet}^2}{\mu_\bullet\Delta t\,r_\bullet^2}.
\]
These remain mean-square / in-probability statements: the STL mean-score noise
$\varepsilon_b=(C^{-1/2}-C^{1/2})Z$ is nonzero whenever $C\neq\Id$, so contraction is \emph{not}
pathwise and the exit probability is \emph{not} zero. The state-dependent constant $V_1=24$ survives
($V_1$ does not contain $\LHs$), so the mean noise persists along any trajectory with $C\neq\Id$.
\item[(b)] \emph{Matched covariance $C_0=\Id$ (i.e.\ $C_0=C_\star$), $m_0\neq0$.} Then
$\varepsilon_b=(C^{-1/2}-C^{1/2})Z=0$ and $\varepsilon_H=\nabla^2\widehat V-A(a)=0$ pathwise, both
covariance updates fix $C_n\equiv\Id$, and the mean update is \emph{deterministic},
$m_{n+1}=(1-\Delta t)m_n$ for both schemes. The iterate is therefore pathwise deterministic and
contracts to $a_\star$ with no stochastic noise---so no exit-probability control is needed---even
though $m_0\neq m_\star$. Thus zero-noise pathwise evolution requires only $C_0=C_\star$, not
$a_0=a_\star$.
\item[(c)] \emph{Quadratic-rescue initialization.} The rescue sets $a_0=a_\star$ exactly
(\Cref{thm:gauss-recovery}), the special case of (b) with $m_0=0$; then $a_n\equiv a_\star$ for all $n$
and the pipeline returns $a_\star$ after the single query. No local phase is needed.
\end{enumerate}
Part~(a) is the local mean-square form of the sticking-the-landing zero-\emph{additive}-floor property
(\Cref{prop:gauss-nofloor}); part~(b) is zero-noise pathwise contraction under matched covariance; and
exactness at $a_\star$ is special to (c).
\end{corollary}

\begin{proof}
\PROOFHOLE
\end{proof}

\subsubsection{Global-to-local entry}
\label{subsec:global-to-local}

The local theorems assume entry into the certified region. The following corollary supplies that
entry from the global stochastic guarantees, completing the stochastic three-stage story with a
single union bound. It uses the two nested regions of \Cref{prop:finite-horizon}: entry is certified
on the inner level $\delta_{\mathrm{in}}$, while the exit time is measured against the outer ball of
radius $r_\bullet$.

\begin{corollary}[Stochastic global-to-local pipeline]
\label{cor:stoch-three-stage}
Fix a scheme $\bullet$, a stochastically admissible outer level $\delta_{\mathrm{out}}$ with radius
$r_\bullet\leq\tfrac1{2\sqrt2}$ satisfying $\{a:\norm{a-a_\star}_{\star,\bullet}\leq r_\bullet\}\subseteq\Ureg_{\delta_{\mathrm{out}}}$,
and stepsize/batch $(\Delta t,B)$ satisfying the hypotheses of \Cref{thm:stoch-local-Riem} (resp.\
\Cref{thm:stoch-local-KL}) at level $\delta_{\mathrm{out}}$. Fix failure budgets
$\epsilon_{\mathrm{ent}},\epsilon'\in(0,1)$ and an inner level $\delta_{\mathrm{in}}\leq\delta_{\mathrm{out}}$
small enough for the deep-entry condition
\[
    \mathfrak b_{\star,\bullet}(\delta_{\mathrm{in}})^2\,\delta_{\mathrm{in}}
    \leq\frac{\epsilon'}2\cdot\frac{\mu_\bullet\Delta t\,r_\bullet^2}{q_\bullet},\qquad q_\bullet=1-\mu_\bullet\Delta t .
\]
Run a stochastic global scheme (\Cref{thm:Riem-Hlip,thm:Riem-selfbound,thm:KL-Hlip,thm:KL-selfbound}
at constant step, or \Cref{thm:decreasing}) for $N_1$ steps until $\E\Delta_{N_1}\leq\epsilon_{\mathrm{ent}}\delta_{\mathrm{in}}$
(achievable when the global floor is below $\epsilon_{\mathrm{ent}}\delta_{\mathrm{in}}$ at constant step,
or for $N_1$ large under the decreasing schedule), then switch to the local scheme with batch
$B\geq c_\bullet N_2\Delta t\,N_\theta^2\LHs^2/(\epsilon'\mu_\bullet r_\bullet^2)$ ($c_{\Riem}=384$,
$c_{\KL}=408$). Write $E_{\mathrm{in}}:=\{a_{N_1}\in\Ureg_{\delta_{\mathrm{in}}}\}$ and the shifted exit
time $\tau^{(N_1)}:=\inf\{n\geq0:\norm{a_{N_1+n}-a_\star}_{\star,\bullet}>r_\bullet\}$. Then for any local
horizon $N_2$,
\[
    \Prob\bigl(E_{\mathrm{in}}\cap\{\tau^{(N_1)}>N_2\}\bigr)\geq 1-\epsilon_{\mathrm{ent}}-\epsilon',
\]
and the entry-restricted non-exit second moment obeys
\[
    \E\bigl[\norm{a_{N_1+N_2}-a_\star}_{\star,\bullet}^2\,;\,E_{\mathrm{in}},\,\tau^{(N_1)}>N_2\bigr]
    \leq(1-\mu_\bullet\Delta t)^{N_2}\,\mathfrak b_{\star,\bullet}(\delta_{\mathrm{in}})^2\delta_{\mathrm{in}}
    +F_\bullet,
\]
with the explicit floors
$F_{\Riem}=\dfrac{384\,\Delta t\,(4+\Gamma)\,N_\theta^2\LHs^2}{B}$ and
$F_{\KL}=\dfrac{408\,\Delta t\,(4+2\Delta t+\Gamma)\,N_\theta^2\LHs^2}{B}$. The restriction to
$E_{\mathrm{in}}$ is necessary: on $E_{\mathrm{in}}^c$ the iterate may lie in the outer ball with
$\norm{a_{N_1}-a_\star}_{\star,\bullet}$ only bounded by $r_\bullet$ (not $r_0$), so the local
contraction does not control that mass.
\end{corollary}

\begin{proof}
\PROOFHOLE
\end{proof}

\begin{remark}[What is and is not proved]
\label{rem:stoch-local-scope}
\Cref{thm:stoch-local-Riem,thm:stoch-local-KL} are \emph{non-exit} (killed-process) statements: they
bound the mean-square deviation on the survival event $\{\tau_\bullet>N\}$, and do not claim a
pathwise or stopped-process bound. They are localized because a bounded local energy sublevel cannot
be pathwise invariant at fixed stepsize: the standing assumptions do not guarantee bounded support for
the STL mean-score noise. Even for a standard Gaussian target, away from $C=\Id$ it equals
$(C^{-1/2}-C^{1/2})Z$ and so has unbounded support---this counterexample alone suffices---so a
projection, clipping, or safeguarding mechanism would be needed for pathwise containment. We do \emph{not} assert almost-sure eventual exit; the present assumptions simply do not
imply infinite-horizon almost-sure containment in a bounded region. Note that the sampled-Hessian
covariance noise $\varepsilon_H$ is, by contrast, bounded a.s.\ ($\alpha_\star\Id\preceq\overline H_n\preceq\beta_\star\Id$),
so it is specifically the mean/score component that obstructs pathwise invariance.
\Cref{prop:finite-horizon} converts the non-exit bound into the operative guarantee: with a deep-entry
condition on $a_0$ and a batch $B=\Omega(N)$, the iterate stays in the band over horizon $N$ with high
probability, and on that event the floor is $O(\Delta t\,A_\bullet N_\theta^2\LHs^2/B)$. The floor
constants ($V_0,V_1$; the aggregate coefficients $64,68$; the exponential-remainder constant $e/2$)
are explicit and conservative, proved in \Cref{lem:local-variance,lem:exp-perturb,lem:local-perturb-Riem,lem:local-perturb-KL};
they are not optimized, and only their orders in $(\Delta t,N_\theta,\LHs,B,A_\bullet)$ enter the
conclusions.
\end{remark}

\subsection{Floor-free convergence via decreasing stepsizes}
\label{sec:decreasing}

At constant stepsize the algorithm converges to a floor. \emph{Under Hessian-Lipschitzness}
(\Cref{assump:global-Hess-Lip}), a common decreasing schedule removes the floor and gives floor-free
$O(1/N)$ convergence for non-Gaussian targets, unconditionally for both schemes. It is stated only in
the Hessian-Lipschitz regime: in the assumption-minimal self-bounding regime the stability threshold
is $\Delta t=\Theta(\kappa^{-2})$ rather than $1/(8\kappa)$, so the schedule
\eqref{eq:decreasing-schedule} does not apply as written and a separate schedule and recursion would
be required.

\begin{theorem}[Decreasing-stepsize floor-free convergence]
\label{thm:decreasing}
Assume \Cref{assump:logconcave-smooth,assump:global-Hess-Lip}, initialize with \eqref{eq:rescue},
and run either \eqref{eq:Riem-STL} or \eqref{eq:KL-STL} with the schedule
\begin{equation}
\label{eq:decreasing-schedule}
    \Delta t_n=\frac{8\kappa}{n+n_0},
    \qquad
    n_0\geq64\kappa^2 .
\end{equation}
Then $\Delta t_n\leq1/(8\kappa)$ for all $n\geq0$, and with $c_{\mathsf S}:=\tfrac74$,
$\Psi_H=N_\theta^2\barLH^2$,
\begin{equation}
\label{eq:decreasing-rate}
    \boxed{\;
    \E\Delta_N\leq\frac{K}{N+n_0},
    \qquad
    K:=\max\Bigl\{\frac{64\,c_{\mathsf S}\,\kappa^3\Psi_H}B,\;n_0\,\Delta(a_0)\Bigr\}.
    \;}
\end{equation}
In particular, with the canonical offset $n_0=\lceil64\kappa^2\rceil$,
\begin{equation}
\label{eq:decreasing-rate-canonical}
    K=\widetilde O\!\Bigl(\frac{\kappa^3N_\theta^2\barLH^2}B+\kappa^2\Delta(a_0)\Bigr);
\end{equation}
for a larger offset $n_0\gg\kappa^2$ the initialization contribution $n_0\Delta(a_0)$ scales with the
chosen $n_0$ rather than $\kappa^2$.
\end{theorem}

\begin{proof}
\PROOFHOLE
\end{proof}

\begin{remark}[Reading the floor-free rate]
\label{rem:decreasing-reading}
The bound \eqref{eq:decreasing-rate} is a genuine $O(1/N)$ decay to zero, not to a floor: the
$\Psi_H$-dependent term is carried inside $K$ and divided by $N+n_0$. For a Gaussian target
$\Psi_H=0$ and $\Delta(a_0)=0$ (\Cref{thm:gauss-recovery}), so $K=0$ and convergence is immediate.
The schedule requires only $\barLH<\infty$ to define $K$; the iteration itself does not use $\LH$.
\end{remark}

\subsection{Gaussian targets: the vanishing floor}
\label{sec:gauss-nofloor}

Every stochastic floor in \Cref{sec:stoch-Riem,sec:stoch-KL} originates from the intrinsic Hessian
fluctuation $\Psi$ entering the one-step estimates \eqref{eq:Riem-onestep-Psi} and
\eqref{eq:KL-onestep-Psi}. The Hessian-Lipschitz floors
(\Cref{thm:Riem-Hlip,thm:KL-Hlip}) display this directly, being proportional to
$\Psi_H=N_\theta^2\barLH^2$; the self-bounding floors (\Cref{thm:Riem-selfbound,thm:KL-selfbound})
no longer show $\Psi$ explicitly only because the coarse universal estimate of
\Cref{lem:Psi-self-bound} has already been substituted. Since $\Psi\equiv0$ for Gaussian targets at
the estimator level, all of them vanish once one returns to the intrinsic one-step; this is the
sticking-the-landing zero-floor property in the intrinsic metric.

\begin{proposition}[Zero intrinsic fluctuation for Gaussian targets]
\label{prop:gauss-nofloor}
If $V(x)=\tfrac12(x-\mu)^\top H(x-\mu)+c$ with $H\succ0$, then $\nabla^2V(X)=H$ for all $X$, so
$A(a)=H$ and $\Psi(a)\equiv0$ for every $a\in\Aspace$; the additive stochastic term in the
\emph{intrinsic} $\Psi$-dependent one-step inequalities of \Cref{sec:stoch-Riem,sec:stoch-KL} (the
Riemannian and KL intrinsic estimates) therefore vanishes \emph{unconditionally} (the coarse
self-bounding recursions, in which the universal bound on $\Psi$ has already been substituted, retain
explicit residual terms and do not visibly vanish).
Consequently, \emph{under the covariance-band and stepsize hypotheses of the corresponding global
one-step theorem}---the band $\tfrac1{2\beta}\Id\preceq C_n\preceq\tfrac1\alpha\Id$ (Riemannian) resp.\
$\tfrac1\beta\Id\preceq C_n\preceq\tfrac1\alpha\Id$ (KL) and the stepsize of \Cref{thm:Riem-Hlip}
resp.\ \Cref{thm:KL-Hlip}, with pathwise band persistence given by the stochastic covariance-band
\Cref{lem:stoch-band} and satisfied from $n=0$ by the quadratic-rescued pipeline
(\Cref{thm:disc-rescue-to-FR})---both schemes contract geometrically with no floor:
\[
    \E\Delta_N\leq(1-c\kappa^{-2})^N\Delta(a_0),
\]
with $c=\tfrac1{24}$ (Riemannian, $B=1$) or $c=\tfrac1{16}$ (KL). Under quadratic rescue $a_0=a_\star$
(\Cref{thm:gauss-recovery}), so $\Delta(a_0)=0$ and the iterates remain exactly at $a_\star$; the
geometric statement is then trivially satisfied but exact. (The factors $c\kappa^{-2}$ are
band-dependent; $\Psi\equiv0$ alone gives the pure descent but not the $\kappa^{-2}$ rate.)
\end{proposition}

\begin{proof}
\PROOFHOLE
\end{proof}

\begin{remark}[Scope of the zero-floor property]
\label{rem:gauss-nofloor-scope}
The zero-floor property is specific to Gaussian targets: it requires the sampled Hessian
$\nabla^2V(X)$ to equal its mean $A(a)$ pathwise, i.e.\ a constant Hessian. For a non-Gaussian
strongly log-concave target, $\Psi$ is strictly positive on any neighborhood where $\nabla^2V$
varies, and the constant-stepsize floor is genuinely nonzero; floor-free convergence then requires the
decreasing schedule of \Cref{thm:decreasing}.
\end{remark}

\subsection{Complexity summary and regime comparison}
\label{sec:regime-comparison}

\Cref{tab:complexity} collects the stochastic complexities. Throughout, the rescue contributes one
deterministic query; the reported iteration counts are those of Phase~I. The four constant-stepsize
global rows and the decreasing-stepsize row are unconditional expectation bounds at single sample
($B=1$); the two local rows are localized non-exit bounds, conditional on certified entry into
$\Ureg_{\delta_\bullet}$, shown with general batch $B$ (minibatching is central to the local survival
result). Here $\widetilde O$ suppresses factors logarithmic in $\kappa$, $N_\theta$,
$\Delta(a_0)/\epsilon$, and $\barLH$.

\begin{table}[t]
\centering
\small
\begin{tabular}{llccc}
\hline
Scheme / regime & controlled qty & floor & contraction & $N$ (to $\epsilon$ above floor) \\
\hline
Riem., self-bounding & $\E\Delta_N$ & $O(\kappa N_\theta)$ & $1-\Theta(\kappa^{-3})$ & $\widetilde O(\kappa^3)$ \\
Riem., Hess.-Lip. & $\E\Delta_N$ & $O(\kappa N_\theta^2\barLH^2)$ & $1-\Theta(\kappa^{-2})$ & $\widetilde O(\kappa^2)$ \\
KL, self-bounding & $\E\Delta_N$ & $O(\kappa N_\theta)$ & $1-\Theta(\kappa^{-3})$ & $\widetilde O(\kappa^3)$ \\
KL, Hess.-Lip. & $\E\Delta_N$ & $O(\kappa N_\theta^2\barLH^2)$ & $1-\Theta(\kappa^{-2})$ & $\widetilde O(\kappa^2)$ \\
\hline
Riem., local (non-exit, $B{\gtrsim}V_1$) & $\E[r_N^2;\tau{>}N]$ & $O(N_\theta^2\LHs^2/B)$ & $1-\Theta((4+\Gamma)^{-2})$ & $\widetilde O((4+\Gamma)^2)$ \\
KL, local (non-exit, $B{\gtrsim}V_1$) & $\E[r_N^2;\tau{>}N]$ & $O(N_\theta^2\LHs^2/B)$ & $1-\Theta((4+\Gamma)^{-2})$ & $\widetilde O((4+\Gamma)^2)$ \\
\hline
Decreasing step & $\E\Delta_N$ & $0$ & --- & $\widetilde O\!\bigl(\tfrac{\kappa^3N_\theta^2\barLH^2+\kappa^2\Delta(a_0)}{\epsilon}\bigr)$ \\
Gaussian target & $\E\Delta_N$ & $0$ & --- & $0$ \\
\hline
\end{tabular}
\caption{Stochastic complexity across regimes; global rows are displayed at $B=1$, local rows retain
general batch $B$. The
iteration counts $N$ are to accuracy $\epsilon$ \emph{above the stated floor} (the constant-stepsize
rows do not converge below their floor). The
self-bounding rows use \Cref{assump:logconcave-smooth}; the Hessian-Lipschitz and decreasing-step
rows additionally use \Cref{assump:global-Hess-Lip}. Minibatching $B=\Theta(\kappa)$ trades the
self-bounding regime to $\widetilde O(\kappa^2)$ rounds at $\widetilde O(\kappa^3)$ total oracle
pairs. The decreasing-step row is at $B=1$ (general $B$ carries a $1/B$ on the stochastic term) and
uses the canonical offset $n_0=\lceil64\kappa^2\rceil$ of \eqref{eq:decreasing-rate-canonical}, under
which the initialization term is $\kappa^2\Delta(a_0)$; it has zero asymptotic floor. The Gaussian
row is exact after the one rescue query. The two \emph{local (non-exit)} rows
(\Cref{sec:stoch-local}) apply once the iterate has entered the certified region
$\Ureg_{\delta_\bullet}$: they control the non-exit (killed) second moment
$\E[\norm{a_N-a_\star}_{\star,\bullet}^2;\tau{>}N]$ (not $\E\Delta_N$; on $\Ureg_{\delta_\bullet}$ these
are related by the coercivity bound $\norm{a-a_\star}_\star^2\leq\mathfrak b(\delta_\bullet)^2\Delta(a)$
of \Cref{cor:band}, the direction used for the entry conversion) at the
target-adaptive local rate (with $\Delta t=\Theta((4+\Gamma)^{-1})$ under the universal window),
paired with the finite-horizon containment of \Cref{prop:finite-horizon} and the entry corollary
\Cref{cor:stoch-three-stage}; their floor is proportional to $\LHs^2\leq\barLH^2$ and its additive part
vanishes for Gaussian targets. The displayed local contraction $1-\Theta((4+\Gamma)^{-2})$ and count
$\widetilde O((4+\Gamma)^2)$ are the large-batch branch $B\gtrsim V_1$; for general $B$ they read
$1-\Theta\bigl((4+\Gamma)^{-2}\min\{1,B/V_1\}\bigr)$ and $\widetilde O\bigl((4+\Gamma)^2\max\{1,V_1/B\}\bigr)$
(\Cref{rem:local-condnum}). In the small-$\Gamma$ regime ($\Gamma=O(1)$, $B\gtrsim V_1$; including
sufficiently near-Gaussian targets) the count is $\widetilde O((4+\Gamma)^2)=O(1)$ in its
$\kappa$-dependence, i.e.\ $O(\log\tfrac1\epsilon)$ (only the squared form follows from the theorem,
though $4+\Gamma=O(1)$ makes $O(4+\Gamma)$ and $O((4+\Gamma)^2)$ coincide). The local $N$-column is a
\emph{round} complexity for contraction of the killed second moment; enforcing the finite-horizon
survival probability of \Cref{prop:finite-horizon} requires the additional horizon-dependent batch
$B=\Omega\bigl(N N_\theta^2\LHs^2/(\epsilon' r_\bullet^2)\bigr)$ and may raise the total oracle
complexity, so these rows are not high-probability oracle-complexity results.}
\label{tab:complexity}
\end{table}

\subparagraph{The $\kappa^3\to\kappa^2$ improvement.}
The single-sample gap between the two constant-stepsize regimes is exactly one factor of $\kappa$.
Under \Cref{assump:logconcave-smooth} alone, the self-bounding estimate
(\Cref{lem:Psi-self-bound}) charges the fluctuation $\Psi$ against the gradient norm, forcing the
stepsize down to $\Delta t=\Theta(\kappa^{-2})$ (through $\rho_B\asymp\kappa$ at $B=1$) and giving
contraction $1-\Theta(\kappa^{-3})$. Under the additional \Cref{assump:global-Hess-Lip}, the
fluctuation is controlled by the target's nonquadraticity directly (\Cref{lem:Psi-Hg}), decoupling
it from the gradient norm; the stepsize rises to $\Delta t=\Theta(\kappa^{-1})$ and the contraction
improves to $1-\Theta(\kappa^{-2})$. The floor changes correspondingly from $O(\kappa N_\theta)$,
which is dimension-linear and target-agnostic, to $O(\kappa N_\theta^2\barLH^2)$, which is
dimension-quadratic but proportional to the squared non-Gaussianity $\barLH^2$ and hence far
smaller for near-Gaussian targets.

\subparagraph{Scope and limitations.}
The stochastic theory presupposes a Price/Hessian oracle (\Cref{def:oracle}); the almost-sure
spectral bound \eqref{eq:pathwise-Hess} it provides is what makes the covariance band forward
invariant pathwise (\Cref{lem:stoch-band}) and the global theorems unconditional. A gradient-only
covariance estimator does not enjoy this bound (\Cref{rem:price}), and the pathwise arguments do not
transfer to it. The Hessian-Lipschitz floor scales as $N_\theta^2$; the whitened per-eigenvalue
refinement that would reduce this to $N_\theta$ is not pursued here. Both the global Riemannian and KL
stochastic rates are unconditional (the KL analysis matches the Riemannian one through the
exponential-retraction identity of \Cref{lem:KL-logtangent}); the local rates are non-exit bounds,
conditional on certified entry and paired with finite-horizon exit control. All constants are stated
explicitly and are those derived in the one-step lemmas; no unspecified absolute constants enter the
boxed bounds.
The stochastic local theory (\Cref{sec:stoch-local}) is deliberately a \emph{non-exit}
(killed-process) statement, not a stopped-process bound: because the STL mean-score noise generally
has unbounded support under Gaussian sampling (the sampled-Hessian covariance noise, by contrast, is
bounded a.s.), and because the standing assumptions do not provide bounded STL mean-score noise,
pathwise invariance of a bounded local sublevel cannot be guaranteed at fixed stepsize without
projection, clipping, or another safeguard (the Gaussian $C\neq\Id$ counterexample $(C^{-1/2}-C^{1/2})Z$
shows no general pathwise-invariance theorem is possible); so the operative guarantee is mean-square
contraction on the survival event plus the
finite-horizon exit-probability bound of \Cref{prop:finite-horizon}. Under a deep-entry condition on
the initialization and a batch $B=\Omega(N)$, a horizon-$N$ run stays in the band with high
probability; \Cref{cor:stoch-three-stage} supplies the high-probability entry from the global phase by
Markov's inequality. We do not claim almost-sure infinite-horizon containment in general; for Gaussian
targets, however, the noise is pathwise zero whenever $C_0=C_\star$, yielding deterministic contraction
and \emph{exact pathwise containment} for any matched-covariance initialization in the certified region
(\Cref{cor:stoch-local-gauss}(b)). Quadratic rescue is the special case $a_0=a_\star$, for which exact
recovery and stationarity ($a_n\equiv a_\star$) occur immediately with no iterative phase. The explicit constants
($V_0,V_1$, the aggregate $64,68$, the exponential remainder $e/2$) are conservative and not
optimized (\Cref{rem:stoch-local-scope}); only their orders enter the conclusions.

\subparagraph{Per-iteration linear algebra.}
The complexities above count Price/Hessian oracle pairs, not dense linear algebra. Each iteration
also performs, in the ambient dimension $N_\theta$: a matrix exponential and a symmetric
factorization for the Riemannian covariance update, or a symmetric positive-definite solve/inverse
for the KL resolvent update, each $O(N_\theta^3)$; $B$ oracle pairs; and $O(N_\theta^2)$ vector work
for the STL score. The one-time rescue is one gradient/Hessian pair and one SPD solve. Thus the
wall-clock cost per iteration is $O(N_\theta^3+B\,c_{\mathrm{oracle}})$, and oracle complexity should
not be read as total arithmetic complexity when $N_\theta$ is large.

\section{Numerical illustrations}
\label{sec:numerics}
\TODO{Confirmatory experiments (one figure + a two-line takeaway each): covariance burn-in scaling;
global localization vs.\ $\kappastar$ on the spiral family (continuous-flow fixed point, not the
discrete map); the fixed-step lower bound (bump train; reproduce the energy ratio $\approx0.273$
across $\kappa\in\{10,20,40,80,160\}$); Gaussian-core exactness of the certified region; the local-rate
frontier (radial-softplus counterexample); and stochastic end-to-end behavior (Monte Carlo on a
non-Gaussian target with analytic energy).}

\section{Discussion and open problems}
\label{sec:discussion}
\TODO{Open problems: the continuous-time global-localization lower bound (the open ledger cell), with a
necessary-conditions checklist for any future construction; the $N_\theta^{p}$, $p\in[1/4,1)$ local-rate
frontier and whether the counterexamples Codazzi/Bochner structure closes it; reducing the variance
floor $N_\theta^2\to N_\theta$ and the global $\kappa$-exponent to $\sqrt\kappa$; removing the squared
$(4+\Gamma)^2$ local factor via KL/Bregman A-stability; and the Gaussian-mixture extension (WFR flow),
where spectral covariance propagation fails unconditionally.}

\appendix

\section{The logarithmic-spiral construction}
\label{app:spiral}
\TODO{Insert the full logarithmic-spiral potential $V_K$ on $\R^2$, the admissible anisotropic
initialization, the global Hessian certificate, the Gaussian shadowing certificate, and the
optimizer-covariance / $\kappastar$ computation supporting \Cref{thm:spiral-sharp,cor:cont-sharp}.}

\section{The bump-train construction}
\label{app:bumptrain}
\TODO{Insert the flat-top curvature bump $\varphi$, the bump train, the geometry of the centers,
regularity and exact condition number, averaged curvature and score at the centers, uniform discrete
shadowing for both covariance maps, and the macroscopic objective-gap lemma supporting
\Cref{thm:disc-global-sharp,cor:disc-global-sharp}.}

\section{Proof of the deterministic KL/Bregman one-step estimate}
\label{app:KL-onestep}

This appendix proves \Cref{prop:KL-onestep} in full, with no imported constants. Throughout, fix
$a=(m,C)$ with $L_n\Id\preceq C\preceq\Lambda\Id$, write $A:=A(a)$, $g:=g(a)$, $R:=C^{-1}-A$,
$\LB=\beta\Lambda\geq1$, and assume $0<\Delta t\leq1/\LB$. Recall $\calE=\calE_1+\calE_2$ with
$\calE_1(a)=-\tfrac12\log\det C+\text{const}$ and $\calE_2(a)=\E_{\rho_a}[V]$, and the Bonnet--Price
identities $\nabla_m\calE_2=\E_{\rho_a}[\nabla V]=-g$, $\nabla_C\calE_2=\tfrac12\E_{\rho_a}[\nabla^2V]=\tfrac12A$.
The descent argument uses only convexity and $\beta$-smoothness of $V$ ($0\preceq\nabla^2V\preceq\beta\Id$),
which hold under \Cref{assump:logconcave-smooth}.

\subsection{Variational characterization and explicit maps}
\label{app:KL-var}

\begin{lemma}
\label{lem:KL-argmin}
The KL/Bregman update \eqref{eq:KL-det} is the forward--backward proximal step
\begin{equation}
\label{eq:KL-argmin}
    a_{n+1}=\arg\min_{a'=(m',C')}\Bigl\{\Delta t\,\inner{\nabla\calE_2(a)}{a'-a}+\Delta t\,\calE_1(a')+\KL(\rho_{a'}\|\rho_a)\Bigr\},
\end{equation}
which linearizes the potential $\calE_2$ explicitly and treats the entropy $\calE_1$ implicitly, and
its minimizer is
\begin{equation}
\label{eq:KL-maps}
    m_{n+1}=m+\Delta t\,C g,
    \qquad
    C_{n+1}=(1+\Delta t)(C^{-1}+\Delta t\,A)^{-1}.
\end{equation}
\end{lemma}

\begin{proof}
\PROOFHOLE
\end{proof}

\subsection{A scalar inequality for the covariance factor}
\label{app:KL-scalar}

\begin{lemma}
\label{lem:scalar-KL}
Let $0<\Delta t\leq1/\LB$, $\mu\in[0,\LB]$, and $b:=\sqrt{(1+\Delta t)/(1+\Delta t\,\mu)}$. Then
\begin{equation}
\label{eq:scalar-KL}
    \mu(b-1)-\log b+\frac{\LB}2(b-1)^2\leq-\frac1{2\Delta t}\bigl(b^2-1-2\log b\bigr).
\end{equation}
\end{lemma}

\begin{proof}
\PROOFHOLE
\end{proof}

\subsection{Forward-KL descent and Wasserstein-gradient lower bound}
\label{app:KL-descent}

\begin{lemma}
\label{lem:KL-onestep-app}
For $0<\Delta t\leq1/\LB$,
\begin{equation}
\label{eq:KL-descent-app}
    \calE(a_{n+1})\leq\calE(a)-\frac1{\Delta t}\KL(\rho_{a_{n+1}}\|\rho_a),
\end{equation}
and, with $\grad^{W_2}\calE(a)=(-g,-R)$ and
$\norm{\grad^{W_2}\calE(a)}_{W_2}^2=\norm g^2+\Tr(RCR)$ (\Cref{lem:W2-PL}),
\begin{equation}
\label{eq:KL-lower}
    \KL(\rho_{a_{n+1}}\|\rho_a)\geq\frac{L_n\Delta t^2}{4(1+\Delta t\,\LB)^2}\,\norm{\grad^{W_2}\calE(a)}_{W_2}^2 .
\end{equation}
\end{lemma}

\begin{proof}
\PROOFHOLE
\end{proof}

Together \Cref{lem:KL-argmin,lem:scalar-KL,lem:KL-onestep-app} prove \Cref{prop:KL-onestep}. Every
constant---the variational maps, the descent margin, and the lower-bound factor
$\tfrac{L_n\Delta t^2}{4(1+\Delta t\LB)^2}$---is derived here; nothing is imported.

\section{Non-Gaussian local modulus: detailed derivation}
\label{app:Kng}

This appendix gives the full auditable proofs of the score identities (\Cref{lem:second-score}) and
the non-Gaussian modulus bound (\Cref{lem:Kng}). We use the whitened notation of \Cref{sec:local}:
$v=C^{-1/2}u$, $S=C^{-1/2}UC^{-1/2}$, $t=\norm\zeta_a$, $X=m+C^{1/2}Z$, $Z\sim\N(0,\Id)$, and
$\bar\Sigma_\delta:=\min\{\Sigma_\delta,\beta_\star-\alpha_\star\}$. Recall the whitened score
generator $\mathcal L$ (Ornstein--Uhlenbeck), which acts as multiplication by $-k$ on the $k$-th
Hermite chaos, so that for $\phi=\sum_{k\geq1}\phi_k$ one has, by Gaussian integration by parts
($-\mathcal L=\nabla^*\nabla$),
$\E\norm{\nabla(-\mathcal L)^{-1}\phi}^2=\E[\phi\,(-\mathcal L)^{-1}\phi]=\sum_{k\geq1}\tfrac1k\norm{\phi_k}_{L^2}^2$.
In particular, when $\phi$ has only chaoses $\geq2$ (as does $\chi_\zeta$),
$\E\norm{\nabla(-\mathcal L)^{-1}\phi}^2\leq\tfrac12\norm\phi_{L^2}^2$; the bound fails for chaos-one
$\phi$ (e.g.\ $\phi=Z_1$ gives equality with constant $1$), but only chaoses $\geq2$ occur below.

\begin{proof}
\PROOFHOLE
\end{proof}

\begin{proof}
\PROOFHOLE
\end{proof}

\bibliographystyle{plainnat}
\bibliography{refs}

\end{document}